\documentclass[11pt,letterpaper]{article}

\usepackage[a4paper,margin=23mm,includeheadfoot,headheight=15pt,headsep=15pt,footskip=22pt]{geometry}
\usepackage{setspace,titlesec,etoolbox,tocloft}
\usepackage[nopatch=footnote]{microtype}
\usepackage{amsmath,amssymb,amsthm,mathtools}
\usepackage{booktabs,longtable,array,tabularx}
\usepackage{graphicx}
\usepackage{float}
\usepackage{xcolor}
\usepackage{enumitem}
\usepackage{fancyhdr}
\usepackage{hyperref}
\usepackage[nameinlink,noabbrev]{cleveref}
\usepackage{bookmark}
\usepackage{listings}
\usepackage{caption,tikz,multicol,xurl}
\usetikzlibrary{arrows.meta,decorations.pathreplacing,patterns,calc}
\usepackage{aliascnt}
\usepackage{needspace}
\usepackage{placeins}

\graphicspath{{core_figures/}}
\allowdisplaybreaks

\hypersetup{
  pdftitle={Projector-Resolved Seed Calculus V},
  pdfauthor={Komorebi P.},
  pdfsubject={Polyphase signal calculus, edge packets, wavelets, and Smith coordinates},
  colorlinks=true,
  linkcolor=black,
  citecolor=black,
  urlcolor=black
}

\pagestyle{fancy}
\fancyhf{}
\fancyhead[L]{\sffamily\fontsize{8}{10}\selectfont\color{muted}KOMOREBI MATHEMATICS / PAPER V}
\fancyhead[R]{\sffamily\fontsize{8}{10}\selectfont\color{muted}TRADITIONAL PAPER}
\fancyfoot[L]{\sffamily\fontsize{8}{10}\selectfont\color{muted}Multiresolution signal calculus}
\fancyfoot[R]{\sffamily\small\color{accent}\thepage}

\newtheorem{theorem}{Theorem}[section]
\theoremstyle{plain}
\newaliascnt{proposition}{theorem}
\newtheorem{proposition}[proposition]{Proposition}
\aliascntresetthe{proposition}
\theoremstyle{plain}
\newaliascnt{lemma}{theorem}
\newtheorem{lemma}[lemma]{Lemma}
\aliascntresetthe{lemma}
\theoremstyle{plain}
\newaliascnt{corollary}{theorem}
\newtheorem{corollary}[corollary]{Corollary}
\aliascntresetthe{corollary}
\theoremstyle{definition}
\newaliascnt{definition}{theorem}
\newtheorem{definition}[definition]{Definition}
\aliascntresetthe{definition}
\theoremstyle{definition}
\newaliascnt{example}{theorem}
\newtheorem{example}[example]{Example}
\aliascntresetthe{example}
\theoremstyle{definition}
\newaliascnt{construction}{theorem}
\newtheorem{construction}[construction]{Construction}
\aliascntresetthe{construction}
\theoremstyle{definition}
\newaliascnt{conjecture}{theorem}
\newtheorem{conjecture}[conjecture]{Conjecture}
\aliascntresetthe{conjecture}
\theoremstyle{remark}
\newaliascnt{remark}{theorem}
\newtheorem{remark}[remark]{Remark}
\aliascntresetthe{remark}
\theoremstyle{remark}
\newaliascnt{boundary}{theorem}
\newtheorem{boundary}[boundary]{Boundary}
\aliascntresetthe{boundary}
\theoremstyle{remark}
\newaliascnt{certificate}{theorem}
\newtheorem{certificate}[certificate]{Numerical Certificate}
\aliascntresetthe{certificate}
\theoremstyle{remark}
\newaliascnt{classical}{theorem}
\newtheorem{classical}[classical]{Classical Input}
\aliascntresetthe{classical}

\newcommand{\cA}{\mathcal A}
\newcommand{\cB}{\mathcal B}
\newcommand{\cE}{\mathcal E}
\newcommand{\cF}{\mathcal F}
\newcommand{\cH}{\mathcal H}
\newcommand{\cP}{\mathcal P}
\newcommand{\cR}{\mathcal R}
\newcommand{\cS}{\mathcal S}
\newcommand{\cT}{\mathcal T}
\newcommand{\bbC}{\mathbb C}
\newcommand{\bbR}{\mathbb R}
\newcommand{\bbZ}{\mathbb Z}
\newcommand{\bbN}{\mathbb N}
\newcommand{\Gam}{\Gamma}
\newcommand{\ii}{\mathrm i}
\newcommand{\RMSE}{\operatorname{RMSE}}
\newcommand{\supp}{\operatorname{supp}}
\newcommand{\diag}{\operatorname{diag}}
\newcommand{\lcm}{\operatorname{lcm}}
\newcommand{\IoU}{\operatorname{IoU}}
\newcommand{\rank}{\operatorname{rank}}
\newcolumntype{L}[1]{>{\raggedright\arraybackslash}p{#1}}
\newcolumntype{Y}{>{\raggedright\arraybackslash}X}

\lstset{
  language=Python,
  basicstyle=\ttfamily\footnotesize,
  breaklines=true,
  frame=single,
  showstringspaces=false,
  columns=fullflexible,
  keepspaces=true,
  numbers=left,
  numberstyle=\tiny,
  numbersep=7pt
}

\title{\textbf{Projector-Resolved Seed Calculus V}\\[0.45em]
\large Polyphase Signal Calculus: Exact Reconstruction, Edge Packets, Wavelets, and Smith Coordinates}
\author{Komorebi P.}
\date{Expanded Lean-checked edition, version 1.1.0\\\small 6 September 2026}


\definecolor{accent}{HTML}{137668}
\definecolor{muted}{HTML}{536568}
\hypersetup{linkcolor=accent,citecolor=accent,urlcolor=accent}
\frenchspacing\clubpenalty=10000\widowpenalty=10000\raggedbottom
\titleformat{\section}{\sffamily\bfseries\large\color{accent}}{\thesection}{.7em}{}
\titleformat{\subsection}{\sffamily\bfseries\normalsize}{\thesubsection}{.7em}{}
\titlespacing*{\section}{0pt}{1.4\baselineskip}{.55\baselineskip}
\pretocmd{\section}{\Needspace{7\baselineskip}}{}{}
\pretocmd{\subsection}{\Needspace{5\baselineskip}}{}{}
\setstretch{1.055}\setlength{\parskip}{.3em}\setlength{\parindent}{1.35em}
\setlength{\emergencystretch}{3em}
\setlength{\cftsecnumwidth}{2.5em}
\setlength{\cftsubsecindent}{2.5em}
\setlength{\cftsubsecnumwidth}{3.4em}
\newaliascnt{standardfact}{theorem}
\newtheorem{standardfact}[standardfact]{Standard Fact}
\aliascntresetthe{standardfact}
\crefname{standardfact}{standard fact}{standard facts}
\crefname{classical}{classical input}{classical inputs}
\crefname{boundary}{boundary}{boundaries}
\crefname{certificate}{numerical certificate}{numerical certificates}
\newcommand{\CC}{\mathbb C}
\newcommand{\RR}{\mathbb R}
\newcommand{\ZZ}{\mathbb Z}
\newcommand{\TT}{\mathbb T}
\newcommand{\Acal}{\mathcal A}
\newcommand{\Ccal}{\mathcal C}
\newcommand{\Ecal}{\mathcal E}
\newcommand{\Fcal}{\mathcal F}
\newcommand{\Hcal}{\mathcal H}
\newcommand{\Kcal}{\mathcal K}
\newcommand{\Lcal}{\mathcal L}
\newcommand{\Pcal}{\mathcal P}
\newcommand{\Qcal}{\mathcal Q}
\newcommand{\Rcal}{\mathcal R}
\newcommand{\Scal}{\mathcal S}
\newcommand{\Tcal}{\mathcal T}
\newcommand{\Wcal}{\mathcal W}
\newcommand{\BV}{\operatorname{BV}}
\newcommand{\dist}{\operatorname{dist}}
\newcommand{\diam}{\operatorname{diam}}
\newcommand{\sgn}{\operatorname{sgn}}
\newcommand{\Per}{\operatorname{Per}}
\newcommand{\id}{\mathrm{id}}
\newcommand{\esssup}{\operatorname*{ess\,sup}}
\newcommand{\one}{\mathbf 1}
\newcommand{\dd}{\,\mathrm d}
\newcommand{\e}{\mathrm e}
\newcommand{\Span}{\operatorname{span}}
\newcommand{\reach}{\operatorname{reach}}
\newcommand{\jump}[1]{\left[#1\right]}

% Operation Komorebi: shared publication profile.
\usepackage{komorebi-monograph}
\renewcommand{\KomorebiTitle}{Multiresolution Signal Calculus}
\renewcommand{\KomorebiSubtitle}{Exact coordinates, singular packets, and stable reconstruction}
\renewcommand{\KomorebiSeries}{Core series V}
\renewcommand{\KomorebiKind}{Modern Monograph}
\renewcommand{\KomorebiEdition}{9 October 2026}
\renewcommand{\KomorebiShortTitle}{Paper V}
\renewcommand{\KomorebiCoverTitle}{Multiresolution\\Signal Calculus}
\renewcommand{\KomorebiCoverBlurb}{Exact coordinates reveal how a signal moves through filters, refinements, and reconstruction. Singular packets keep representation, approximation, and recovery distinct.}
\renewcommand{\KomorebiTitleSize}{35}
\renewcommand{\KomorebiTitleLeading}{40}
\renewcommand{\KomorebiProfile}{monograph}
\renewcommand{\KomorebiOrientation}{portrait}
\KomorebiSetup
% Preserve legacy body tokens while rendering bold vectors and unbraced binary operators with Unicode math.
\providecommand{\bm}[1]{\symbfit{#1}}
\let\KomorebiPrimitiveMathbin\mathbin
\renewcommand{\mathbin}[1]{\KomorebiPrimitiveMathbin{#1}}
\let\KomorebiPrimitiveWidehat\widehat
\renewcommand{\widehat}[1]{\KomorebiPrimitiveWidehat{#1}}
\crefname{proposition}{proposition}{propositions}
\crefname{lemma}{lemma}{lemmas}
\crefname{corollary}{corollary}{corollaries}
\crefname{definition}{definition}{definitions}
\crefname{example}{example}{examples}
\crefname{construction}{construction}{constructions}
\crefname{conjecture}{conjecture}{conjectures}
\crefname{remark}{remark}{remarks}
\crefname{boundary}{boundary}{boundaries}
\crefname{certificate}{numerical certificate}{numerical certificates}
\crefname{classical}{classical input}{classical inputs}
\input{atlas-definitions.tex}
\begin{document}
\KomorebiCover
\KomorebiOpeningNote{\textbf{Part I.} The twenty-result formal core develops exact polyphase reconstruction, filter transfers, Hilbert packet operators, finite adaptive trees, edge error bounds, and Smith coordinates.

\medskip
\textbf{Part II.} The complete extended development treats lattice coordinates, stable banks, singular packets, separated jump recovery, geometric error bounds, and multiscale localization.

\medskip
\textbf{Visual atlas.} Forty worked plates close the volume. The separate visual guide provides the same examples, and the Lean guide gives the precise released formal contract.}

\section*{Reading this edition}\addcontentsline{toc}{section}{Reading this edition}
The two branches of Paper V answer related questions at different levels. Part I preserves the twenty results and twenty proofs of the later v1.1.0 manuscript. Part II preserves the broader standalone development with its sixty-three proofs and E-prefixed section numbers. Their overlap is intentional. The combined eighty-three proofs are a record of the two complete branches, rather than eighty-three independent or novel results.

For a first reading, follow Part I Sections 1--7. These sections establish the coordinate maps, filter transfers, exact reconstruction, and the separation of packet error from smooth approximation error. Section 11 develops finite adaptive selection. Sections 13--14 give the cyclic Smith realization. Continue with Part II Sections E2--E6 for lattice and stability generalizations, E7--E11 for singular packets and recovery, and E12--E14 for nonlinear and Cayley coordinates. Each visual plate identifies the corresponding manuscript result.

The inherited September 2026 release records 225 Lean theorem targets across 23 proof modules supporting all twenty Part I results. This scope does not extend automatically to Part II. The E-prefixed development has prose proofs and retains the mathematical corrections recorded in the prior edition's exact change ledger. No fresh Lean compilation is claimed for this monograph. The current computation review reran twenty exact-model tests with 305,933 counted assertions and sixteen verifier controls. The current companion suite passed thirty-three tests, with twenty-seven overlapping tests from the earlier snapshot also checked. The extended suite passed 11,052 checks in thirteen groups. Fifty-two saved floating-point metrics and eighteen discrete leaves reproduced exactly, and all eighteen companion figure pairs were regenerated. The visual examples passed their own symbolic checks and controls. These checks have finite scopes and are recorded separately in the accompanying verification files. Historical counts in Part I's reproducibility appendix keep their September provenance.

The two broad edge recovery conjectures remain open. Part II proves recovery for a declared one dimensional separated jump class under its sampling, contrast, and noise assumptions. These results do not supply universal image segmentation or infinite-tree consistency. Numerical illustrations keep the finite interpretation stated in their captions.
\clearpage
\tableofcontents
\clearpage
\listoffigures
\listoftables
\clearpage
\pagenumbering{arabic}
\part*{Part I: The released formal core}\addcontentsline{toc}{part}{Part I: The released formal core}
\begin{readerroadmap}[title={Coordinates first, decisions second}]
A residue decomposition changes the coordinates of a signal without changing its information. Filtering then acts between those coordinates. Approximation enters when an operator compresses a component or when a decision discards coefficients. The formal core follows these operations in that order.
\end{readerroadmap}



\begin{abstract}
This paper develops the signal realization of the projector-resolved seed calculus. The impulse $\delta_0$ is taken as a primitive signal seed and its shift orbit supplies the address system. Residue projectors split a sequence into exact polyphase channels, nested divisibility gives exact refinement, and every finite impulse-response operator becomes a matrix of packet transfers. The standard perfect-reconstruction condition for a filter bank is therefore an analysis--synthesis identity in the common operator language of the series. Haar analysis and every finite-depth wavelet-packet transform are treated on complex $\ell^2$, together with concrete finite-signal models and a predict--update lifting step. The edge problem is separated into two statements. When a discontinuity packet is known and retained exactly, approximation error is exactly the approximation error of the remaining smooth packet. No spectral approximation of the jump is performed. When the packet is estimated, a symmetric-difference bound isolates mask error, amplitude error, and smooth-background error. Numerical studies compare global Fourier truncation, exact known-packet reconstruction, and a deliberately narrow positive-contrast estimator in one and two dimensions, including curved boundaries, corners, and multiple amplitudes. The same cyclic phase action is then placed on the Smith reflection disk. Its pullback through the Cayley transform is a M\"obius action on impedance, and the two-phase case is precisely impedance--admittance inversion. Nonlinear coefficient products obey cyclic packet convolution. The preserved September 2026 release integrates a Lean formalization of all twenty theorem-like statements under the hypotheses stated below: 225 supporting theorem targets in 23 modules. Eighteen regenerated figures, machine-readable numerical results, the 33-test original companion suite, and 36 additional exact-model and verifier tests accompany the paper. Exact signal theorems, limited numerical estimators, classical constructions, and open recovery problems are kept distinct throughout.
\end{abstract}

\begin{interpretation}[title={Scope and claim boundary}]
 The paper proves exact polyphase reconstruction, exact finite filter-transfer identities, perfect reconstruction for the stated filter banks, the conditional known-edge theorem, mask and amplitude perturbation bounds, nonlinear residue-convolution laws, and the Smith-coordinate phase action. It does not claim that every discontinuity can be detected automatically, that every form of Gibbs ringing has been eliminated, or that the simple two-class estimator supplied with the release is a universal segmentation method. The estimator is included as one reproducible finite example.  For every finite packet refinement tree, the paper now proves existence of a globally optimal complete pruning, an exact dynamic-programming selector for additive costs, and a deterministic margin-stability certificate. Wavelets, filter banks, the Cayley transform, and the classical Smith grid are treated as classical inputs organized in the common seed--packet--refinement language.
\end{interpretation}

\clearpage
\noindent\textbf{Verification convention.} ``Lean-checked'' refers to the explicitly scoped theorem statements and their supporting declarations, accepted by the pinned Lean kernel. It does not certify the English text automatically, the floating-point Python implementation, physical modelling choices, or the two conjectures. The manuscript-to-declaration correspondence has been reviewed manually. Appendix~\ref{sec:formalization} records it result by result. The September 2026 release records a rerun of the numerical experiments. The current final-edition checks are distinguished in the opening reading note.



\section{The signal realization of the seed calculus}\label{core:section:01}

Let $\ell^2(\bbZ;\bbC)$, abbreviated $\ell^2(\bbZ)$, be the Hilbert space of square-summable complex discrete signals. The inner product is conjugate-linear in its first argument. A superscript $*$ on a bounded operator denotes its Hilbert adjoint. Write $\delta_n$ for the unit impulse at address $n$ and let the bilateral shift be
\[
(Sx)[n]=x[n-1].
\]
Then
\[
S^n\delta_0=\delta_n.
\]
The primitive signal seed is therefore $s=\delta_0$, the address operators are the shift words $S^n$, and the completion is the $\ell^2$ norm.

\begin{theorem}[Impulse seed theorem]\label{result:01}
The closed shift-generated hull of $\delta_0$ is all of $\ell^2(\bbZ)$:
\[
\boxed{
\overline{\operatorname{span}\{S^n\delta_0:n\in\bbZ\}}^{\,\ell^2}
=\ell^2(\bbZ).
}
\]
Every finitely supported signal has the exact finite expression
\[
x=\sum_{n\in\bbZ}x[n]S^n\delta_0.
\]
\end{theorem}

\begin{proof}
The shift orbit is the standard orthonormal basis $\{\delta_n\}_{n\in\bbZ}$. For $x\in\ell^2$, the finite truncations
\[
x_F=\sum_{n\in F}x[n]\delta_n,\qquad F\subset\bbZ\text{ finite},
\]
satisfy $\|x-x_F\|_2^2=\sum_{n\notin F}|x[n]|^2\to0$ as the finite sets exhaust $\bbZ$. Thus the unordered Hilbert-space sum $\sum_nx[n]\delta_n$ converges to $x$, and the finite-support statement is exact. The shift is a surjective isometry because $n\mapsto n-1$ is a bijection of $\bbZ$, with inverse shift $n\mapsto n+1$. This proves the closed-span assertion.
\end{proof}

This realization fixes the language used below:
\[
\begin{gathered}
\text{impulse seed}\longrightarrow\text{shift address}\longrightarrow\text{packet projection}\longrightarrow\text{filter transfer}\\
\longrightarrow\text{refinement}\longrightarrow\text{reconstruction}.
\end{gathered}
\]

\begin{remark}[Several completions]\label{core:record:02}
The same algebraic seed description applies to finite sequences, $\ell^1$, $\ell^2$, periodic sequences, tempered distributions, and multidimensional lattices. The topology decides which infinite reconstructions are legitimate. Finite-dimensional experiments in this paper use ordinary Euclidean norms; theorem statements specify the ambient norm when it matters.
\end{remark}

\section{Residue projectors and exact polyphase coordinates}\label{core:section:02}

\begin{keyidea}[title={A packet has two useful representations}]
The masked packet keeps its original sample addresses and sets the other phases to zero. The decimated channel stores only the retained addresses, now indexed by the coarse coordinate. They carry the same values, but their domains differ. Refinement sums belong on the original lattice. Interleaving converts the channels back to that lattice.
\end{keyidea}

Fix an integer $N\ge1$. For $0\le r<N$, define
\[
(P_{N,r}x)[n]
=
\mathbf 1_{n\equiv r\, (\mathrm{mod}\,N)}x[n].
\]
The corresponding decimated channel is
\[
x_r[q]=x[Nq+r].
\]

\begin{theorem}[Polyphase projector theorem]\label{result:02}
The maps $P_{N,r}$ are bounded orthogonal projectors on $\ell^2$, with operator norm at most one, and satisfy
\[
P_{N,r}P_{N,s}=\delta_{rs}P_{N,r},
\qquad
\sum_{r=0}^{N-1}P_{N,r}=I.
\]
Consequently,
\[
\boxed{x=\sum_{r=0}^{N-1}P_{N,r}x}
\]
with no approximation.
\end{theorem}

\begin{proof}
Each integer belongs to exactly one residue class modulo $N$. Multiplying the two indicator functions gives zero for distinct classes and the same indicator for equal classes. Summing all residue indicators gives one at every address. Moreover,
\[
\|P_{N,r}x\|_2^2=\sum_{n\equiv r\ (\mathrm{mod}\,N)}|x[n]|^2\le\|x\|_2^2,
\qquad
\langle P_{N,r}x,y\rangle=\langle x,P_{N,r}y\rangle.
\]
The mask is real-valued, so the inner-product identity follows term by term. Hence each projector is self-adjoint, and the ranges of different residue projectors are orthogonal.
\end{proof}

Define the analysis and synthesis maps
\[
\cA_Nx=(x_0,\ldots,x_{N-1}),
\qquad
\cS_N(x_0,\ldots,x_{N-1})[Nq+r]=x_r[q].
\]

\begin{corollary}[Analysis--synthesis identity]\label{result:03}
\[
\boxed{\cS_N\cA_N=I.}
\]
On the product of the channel spaces one also has $\cA_N\cS_N=I$. Equip this finite product with the Hilbert norm
\[
\|(x_0,\ldots,x_{N-1})\|_{\mathrm{ch}}^2
=\sum_{r=0}^{N-1}\|x_r\|_2^2.
\]
Then $\cA_N$ is a linear isometric equivalence and
\[
\|x\|_2^2=\sum_{r=0}^{N-1}\|x_r\|_2^2.
\]
\end{corollary}
\begin{proof}
Euclidean division gives a unique pair $(r,q)$ with $0\le r<N$ and $n=Nq+r$, including negative $n$. The address maps $n\mapsto(n\bmod N,n\mathbin{\mathrm{div}}N)$ and $(r,q)\mapsto Nq+r$ are inverse. Analysis and synthesis are reindexing by this bijection. Reindexing the nonnegative sum $\sum_n|x[n]|^2$ proves the norm identity and shows that both maps preserve square summability. In Lean the target is $\ell^2(\{0,\ldots,N-1\}\times\bbZ)$, which realizes exactly this Hilbert norm.
\end{proof}

\begin{figure}[htbp]
\centering
\includegraphics[width=0.94\textwidth]{01_polyphase_analysis.pdf}
\caption{Four-phase analysis of a finite signal. The channels are decimated coordinate packets; interleaving them gives the original signal exactly.}
\label{fig:polyphase-analysis}
\end{figure}

\subsection{Nested refinement}

Suppose $N,M\ge1$ and $N\mid M$. Every coarse phase has a disjoint family of fine descendants:
\[
\operatorname{Child}_{M}(N,r)
=
\{s:0\le s<M,\ s\equiv r\pmod N\}.
\]

\begin{theorem}[Exact nested refinement]\label{result:04}
For $N,M\ge1$ and $N\mid M$,
\[
\boxed{
P_{N,r}
=
\sum_{\substack{0\le s<M\\s\equiv r\, (\mathrm{mod}\,N)}}P_{M,s}.
}
\]
\end{theorem}

\begin{proof}
An integer $n$ satisfies $n\equiv r\pmod N$ exactly when its unique residue $s=n\bmod M$ lies among the listed descendants. The fine indicators are disjoint, so their sum is the coarse indicator.
\end{proof}

This identity is the signal form of the refinement law introduced in Paper II. It is not merely a denser sampling rule. The coarse masked packet, represented on the original integer lattice, is exactly the sum of its fine masked descendants. Decimated channels at different moduli must first be placed in these common coordinates before being added.

\section{Linear time-invariant operators as transfer matrices}\label{core:section:03}

Let $h$ be a finite impulse response and
\[
y=h*x,
\qquad
y[n]=\sum_{q\in\bbZ}h[q]x[n-q].
\]
Resolve both input and output into $N$ phases. For output phase $k$ and input phase $s$, define the transfer kernel
\[
\boxed{
h_{ks}[m]=h[Nm+k-s].
}
\]
Out-of-range taps are understood to be zero.

\begin{theorem}[Exact FIR polyphase transfer law]\label{result:05}
The output channels satisfy
\[
\boxed{
y_k[n]
=
\sum_{s=0}^{N-1}\sum_{m\in\bbZ}
h_{ks}[m]x_s[n-m].
}
\]
For the numerical channel transform, take $z\in\bbC\setminus\{0\}$ and assume
\[
\sum_{n\in\bbZ}|x_s[n]z^{-n}|<\infty\quad(0\le s<N).
\]
Define $X_s(z)=\sum_nx_s[n]z^{-n}$. The output series then also converge absolutely, and
\[
\bm Y(z)=\bm H^{(N)}(z)\bm X(z),
\qquad
H^{(N)}_{ks}(z)=\sum_m h[Nm+k-s]z^{-m}.
\]
\end{theorem}

\begin{proof}
Evaluate $y$ at address $Nn+k$ and partition every input index by its residue $s$:
\[
\begin{aligned}
y[Nn+k]
&=\sum_qh[q]x[Nn+k-q]\\
&=\sum_{s=0}^{N-1}\sum_m
h[Nm+k-s]x[N(n-m)+s].
\end{aligned}
\]
For a fixed output phase $k$, the change of tap address $q=Nm+k-s$ is a bijection between $\bbZ$ and the pairs $(s,m)$ with $0\le s<N$: choose $s\equiv k-q\pmod N$ and then $m=(q-k+s)/N$. Thus this is a reindexing of the finite filter support, valid for every bilateral input sequence, without a convergence assumption on $x$.

The last factor is $x_s[n-m]$, giving the time-domain formula. Under the additional transform hypothesis, each delayed channel has transform
\[
\sum_nx_s[n-m]z^{-n}
=z^{-m}\sum_jx_s[j]z^{-j}.
\]
Absolute summability permits this bijective reindexing, and only finitely many delayed channels are combined. Taking their finite linear combination gives the matrix formula. The time-domain theorem therefore remains valid even when no numerical bilateral transform exists.
\end{proof}

\begin{figure}[htbp]
\centering
\includegraphics[width=0.80\textwidth]{02_filter_transfer_matrix.pdf}
\caption{The exact transfer matrix for a seven-tap FIR filter in four phases. Every cell records the delays and coefficients carrying one input packet into one output packet.}
\label{fig:filter-matrix}
\end{figure}

\begin{remark}[Operator-block form]\label{core:record:07}
If $H$ denotes convolution by $h$, the same identity is
\[
H=\sum_{k,s}P_{N,k}HP_{N,s}.
\]
Thus the familiar polyphase matrix is exactly the transfer-block matrix of the seed calculus.
\end{remark}

\section{Analysis banks, synthesis banks, and perfect reconstruction}\label{core:section:04}

\begin{interpretation}[title={Read the delay in the coordinate system used}]
A step in one decimated channel crosses a full block of original samples. A one-sample shift also permutes the residue labels. At two phases, these operations have different matrices. The following calculation keeps the original sample delay visible instead of hiding it in a channel convention.
\end{interpretation}
\begin{figure}[htbp]\centering
\includegraphics[width=.91\linewidth]{monograph_figures/delay_coordinates.pdf}
\caption{Exact address comparison for a two-phase split. A delay of one original sample sends the odd channel into the even channel with a coarse delay, and the even channel into the odd channel without a coarse delay. A one-step channel delay moves both original residues by two samples. All plotted values are integer sample addresses.}\label{monograph:fig:delay}
\end{figure}

Work with finite impulse-response banks. Let $\mathcal L=\bbC[z,z^{-1}]$ be the Laurent-polynomial algebra, with the monomial $z^{-m}$ acting on a bilateral sequence by $x[n]\mapsto x[n-m]$. Let $\bm E\in\operatorname{Mat}_{N\times N}(\mathcal L)$ and $\bm R\in\operatorname{Mat}_{N\times N}(\mathcal L)$ be analysis and synthesis matrices. Their composition acts on all bilateral channel sequences by finite convolutions. At a nonzero $z$ where the input transforms are absolutely convergent, this gives
\[
\widehat{\bm X}(z)=\bm R(z)\bm E(z)\bm X(z).
\]

\begin{theorem}[Perfect-reconstruction criterion]\label{result:06}
For $d\in\bbZ$, such a bank reconstructs every input up to a delay of $d$ \emph{channel indices} exactly when the following Laurent-polynomial identity holds:
\[
\boxed{
\bm R(z)\bm E(z)=z^{-d}I_N.
}
\]
If the right side is $c z^{-d}I_N$ with $c\ne0$, reconstruction holds up to gain $c$ and channel delay $d$. After interleaving, this is a delay of $Nd$ original samples. A negative delay denotes an advance.
\end{theorem}

\begin{proof}
Finite convolution represents Laurent-polynomial multiplication: composing shifts adds their delays, and distributivity handles finite sums. Thus the matrix identity implies the asserted action. Conversely, let $\bm B=\bm R\bm E-cz^{-d}I_N$ and suppose its action vanishes on every channel sequence. Insert an impulse at channel $s$, address zero, with every other channel zero. Output channel $k$ reads the coefficients of the Laurent polynomial $B_{ks}$ at all addresses. All of these coefficients vanish, so $B_{ks}=0$ for every $k,s$. The representation is faithful, proving necessity as a polynomial identity rather than merely equality at one evaluation point. The displayed unit-gain case is $c=1$.
\end{proof}

\paragraph{Original-sample delay.}
For an arbitrary original-sample delay $D\in\bbZ$, the output phase $k$ has
\[
y_k[n]=x_{r_k}[n+q_k],\qquad
r_k=(k-D)\bmod N,\qquad q_k=(k-D)\mathbin{\mathrm{div}}N.
\]
Its channel matrix has the sole nonzero entry $z^{q_k}$ in row $k$, column $r_k$. In particular, a one-sample delay at $N=2$ has matrix
\[
\begin{pmatrix}0&z^{-1}\\1&0\end{pmatrix},
\]
whereas $z^{-1}I_2$ delays the original signal by two samples. The general address identity is valid for negative sample addresses as well.

\subsection{The Haar bank}

The Haar analysis matrix is
\[
\bm E_{\mathrm{H}}
=
\frac1{\sqrt2}
\begin{pmatrix}
1&1\\
1&-1
\end{pmatrix}.
\]
It is orthogonal, so $\bm R_{\mathrm{H}}=\bm E_{\mathrm{H}}^{\mathsf T}$ and
\[
\bm R_{\mathrm{H}}\bm E_{\mathrm{H}}=I_2.
\]
The channels are
\[
a[n]=\frac{x[2n]+x[2n+1]}{\sqrt2},
\qquad
d[n]=\frac{x[2n]-x[2n+1]}{\sqrt2}.
\]

\begin{corollary}[Haar energy identity]\label{result:07}
\[
\boxed{
\|x\|_2^2=\|a\|_2^2+\|d\|_2^2.
}
\]
\end{corollary}

\begin{proof}
For any $u,v\in\bbC$, the parallelogram identity gives
\[
\left|\frac{u+v}{\sqrt2}\right|^2+
\left|\frac{u-v}{\sqrt2}\right|^2=|u|^2+|v|^2.
\]
Apply this to $u=x[2n]$, $v=x[2n+1]$ and sum over $n\in\bbZ$. Each summand is nonnegative; the right side is $\|x\|_2^2$ by the two-channel address equivalence. This also proves that the literal low and detail channels belong to $\ell^2$. The real Haar matrix is unitary over $\bbC$, so the finite-signal statement is the same calculation with a finite sum.
\end{proof}

\begin{figure}[htbp]
\centering
\includegraphics[width=0.92\textwidth]{03_haar_edge_packets.pdf}
\caption{A piecewise-smooth signal and its first Haar low/detail packets. The low channel records local averages. The detail channel records adjacent-sample differences, including pairs that cross a jump.}
\label{fig:haar-packets}
\end{figure}

\subsection{Predict--update lifting}

Split $x$ into even and odd channels $e,o$. Define
\[
d=o-e,
\qquad
s=e+\frac12d.
\]
The inverse is
\[
e=s-\frac12d,
\qquad
o=d+e.
\]

\begin{theorem}[Exact lifting reconstruction]\label{result:08}
The predict--update map $(e,o)\mapsto(s,d)$ is invertible with determinant one and the displayed inverse. Consequently the associated finite transform reconstructs every even-length signal exactly.
\end{theorem}

\begin{proof}
The forward matrix is
\[
\begin{pmatrix}s\\d\end{pmatrix}
=
\begin{pmatrix}1/2&1/2\\-1&1\end{pmatrix}
\begin{pmatrix}e\\o\end{pmatrix},
\]
whose determinant is one. Direct substitution gives the inverse formulas.
\end{proof}

\section{Wavelet packets as addressed refinements}\label{core:section:05}

On $\ell^2(\bbZ;\bbC)$, apply the Haar analysis to both children at each stage. For a finite input, assume its length is divisible by $2^\ell$ when a depth-$\ell$ transform is used. A binary word $\alpha\in\{0,1\}^{\ell}$ addresses a depth-$\ell$ packet, denoted $W_{\ell,\alpha}x$. Appending $0$ or $1$ selects the low or high child.

\begin{theorem}[Wavelet-packet reconstruction and energy]\label{result:09}
At every finite depth $\ell$,
\[
\boxed{
x=\sum_{\alpha\in\{0,1\}^{\ell}}W_{\ell,\alpha}^{*}W_{\ell,\alpha}x,
}
\]
and
\[
\boxed{
\|x\|_2^2
=
\sum_{\alpha\in\{0,1\}^{\ell}}
\|W_{\ell,\alpha}x\|_2^2.
}
\]
\end{theorem}

\begin{proof}
Write $H_0,H_1$ for the literal low and detail operators. The preceding energy identity implies $\|H_bx\|_2\le\|x\|_2$, so both have bounded Hilbert adjoints. At depth zero the only packet operator is the identity. If the energy identity holds at depth $\ell$, apply the one-level identity to each $W_{\ell,\alpha}x$ and sum; this proves it at depth $\ell+1$.

Set $A=\sum_\alpha W_{\ell,\alpha}^*W_{\ell,\alpha}$. This is a finite sum of bounded operators. The energy result says $\langle x,Ax\rangle=\langle x,x\rangle$ for every $x$. Complex polarization, applied also to $x+y$ and $x+\ii y$, gives $\langle y,(A-I)x\rangle=0$ for all $x,y$, hence $A=I$. Applying the operator identity to $x$ proves synthesis. This argument uses the actual Hilbert adjoints of the decimation operators, not a separately postulated inverse symbol.
\end{proof}

\begin{figure}[htbp]
\centering
\includegraphics[width=0.92\textwidth]{04_wavelet_packet_energy.pdf}
\caption{Energy of the sixteen depth-four Haar packet addresses for the test signal. The packet address is an operator word in the common seed grammar.}
\label{fig:wavelet-energy}
\end{figure}

\begin{classical}[Wavelet boundary]\label{core:record:12}
The existence and perfect-reconstruction theory of wavelets and filter banks are classical. Their role here is structural: they provide a fully developed realization in which seed generation, packet addressing, refinement, transfer blocks, and exact synthesis can all be read literally in the common operator language.
\end{classical}

\section{Known edge packets and spectral reconstruction}\label{core:section:06}

\begin{keyidea}[title={Retention changes the approximation target}]
The retained packet is supplied as part of the model. It passes through synthesis exactly. The approximation operator acts on the smooth remainder, so its error can be studied without forcing a global smooth basis to reproduce the jump. Finding that packet from observations is a separate recovery problem.
\end{keyidea}
\begin{figure}[htbp]\centering
\includegraphics[width=.95\linewidth]{monograph_figures/retained_packet_cancellation.pdf}
\caption{An exact rational eight-coordinate example of retained packet cancellation. The supplied jump packet is kept, while the linear operator averages adjacent pairs in the smooth component. The two error vectors agree coordinate by coordinate. This model illustrates cancellation and does not test packet detection or a Fourier convergence rate.}\label{monograph:fig:cancellation}
\end{figure}


The central distinction is between approximating a jump globally and retaining it as its own packet.

Let
\[
f=s+j
\]
where $s$ belongs to a smooth approximation class and $j$ is a known discontinuity packet. Let $Q_m$ be any approximation operator applied only to $s$, and define
\[
\widetilde f_m=Q_ms+j.
\]

\begin{theorem}[Exact known-edge theorem]\label{result:10}
For every norm for which the expressions are defined,
\[
\boxed{
f-\widetilde f_m=s-Q_ms,
\qquad
\|f-\widetilde f_m\|=\|s-Q_ms\|.
}
\]
No approximation of $j$ occurs.
\end{theorem}

\begin{proof}
Subtract:
\[
f-\widetilde f_m=(s+j)-(Q_ms+j)=s-Q_ms.
\]
Taking the norm gives the second identity.
\end{proof}

\begin{corollary}[Several exact edge packets]\label{result:11}
If
\[
f=s+\sum_{q=1}^{M}j_q
\]
and every $j_q$ is retained exactly, then
\[
\left\|f-\left(Q_ms+\sum_qj_q\right)\right\|
=\|s-Q_ms\|.
\]
Curved edges, corners, and different jump amplitudes do not alter the identity.
\end{corollary}
\begin{proof}
The finite sum $j=\sum_{q=1}^M j_q$ is a single retained packet. Applying the preceding cancellation identity gives the result; the empty family $M=0$ is included with sum zero.
\end{proof}

\begin{boundary}[What the theorem does and does not say]\label{core:record:15}
The theorem explains why global ringing caused by spectrally approximating a known jump is absent from the reconstructed jump packet. It does not discover the packet. If the packet is wrong, its error is reconstructed exactly as well. Nor does the theorem prevent unrelated truncation errors in the smooth packet, domain-boundary errors, noise, or model mismatch.
\end{boundary}

\section{Mask, amplitude, and smooth-background error}\label{core:section:07}

Let $(X,\Sigma,\mu)$ be a measure space, with measurable masks $\Omega,\widehat\Omega\subseteq X$ and scalar amplitudes $a,\widehat a\in\bbC$ (real amplitudes are a special case). Suppose a scalar jump packet is
\[
j=a\mathbf 1_{\Omega},
\]
and an estimated packet is
\[
\widehat j=\widehat a\mathbf 1_{\widehat\Omega}.
\]

\begin{theorem}[Packet perturbation bound]\label{result:12}
For $1\le p<\infty$, interpreted first as an inequality of extended $L^p$ seminorms,
\[
\boxed{
\|j-\widehat j\|_{L^p}
\le
|a-\widehat a|\,\mu(\Omega)^{1/p}
+
|\widehat a|\,\mu(\Omega\triangle\widehat\Omega)^{1/p}.
}
\]
If $\mu(\Omega)$ and $\mu(\Omega\triangle\widehat\Omega)$ are finite, all terms in this displayed jump bound are finite and it is an ordinary real-valued $L^p$ inequality. For the pointwise supremum, and consequently also the essential $L^\infty$ seminorm,
\[
\|j-\widehat j\|_{\infty}
\le |a-\widehat a|+|\widehat a|\,
\mathbf 1_{\Omega\triangle\widehat\Omega\ne\varnothing}.
\]
\end{theorem}

\begin{proof}
Write
\[
a\mathbf 1_{\Omega}-\widehat a\mathbf 1_{\widehat\Omega}
=(a-\widehat a)\mathbf 1_{\Omega}
+\widehat a(\mathbf 1_{\Omega}-\mathbf 1_{\widehat\Omega}).
\]
Apply the triangle inequality and use
\[
\|\mathbf 1_{\Omega}-\mathbf 1_{\widehat\Omega}\|_{L^p}
=\mu(\Omega\triangle\widehat\Omega)^{1/p}.
\]
For the supremum estimate, $|\mathbf1_\Omega-\mathbf1_{\widehat\Omega}|\le1$ pointwise, and this function vanishes everywhere when the symmetric difference is empty. The same pointwise bound also bounds the essential supremum. Nonempty null symmetric differences need not contribute to the essential norm; the stated bound remains valid in that case and need not be sharp. Measurability justifies the $L^p$ triangle inequality and the indicator norm formula.
\end{proof}

If $s,\widehat s$ are measurable smooth-background models, the same triangle inequality gives, for finite $p$,
\[
\boxed{
\|f-(\widehat s+\widehat j)\|_p
\le
\|s-\widehat s\|_p
+|a-\widehat a|\mu(\Omega)^{1/p}
+|\widehat a|\mu(\Omega\triangle\widehat\Omega)^{1/p}.
}
\]
This is the correct separation of errors. Detection error is not hidden inside a generic reconstruction score.

\begin{conjecture}[Stable edge-packet recovery]\label{core:record:17}
For specified piecewise-smooth classes, contrast bounds, separation conditions, and noise models, there should exist a multiscale packet estimator whose symmetric-difference error and amplitude error converge at rates controlled independently of the smooth approximation rate. A theorem of this kind would turn the conditional known-edge result into a stable data-driven reconstruction theorem. No universal version is asserted here.
\end{conjecture}

\Needspace{12\baselineskip}
\section{One-dimensional numerical study}\label{core:section:08}

The one-dimensional signal contains a smooth non-bandlimited background and a positive rectangular packet. Three reconstructions are compared:

\begin{enumerate}[label=(\roman*)]
\item global Fourier truncation of the complete signal;
\item Fourier truncation of the known smooth packet plus the exact jump packet;
\item a two-class estimated packet followed by harmonic background filling, spectral smoothing, and jump-amplitude estimation.
\end{enumerate}

\begin{figure}[htbp]
\centering
\includegraphics[width=0.98\textwidth]{05_edge_reconstruction_1d.pdf}
\caption{One-dimensional reconstruction and absolute error. The known packet removes the edge from the global spectral approximation. The estimated packet is visibly better than the global reconstruction in this limited test but retains smooth-background estimation error.}
\label{fig:edge-1d}
\end{figure}

\begin{certificate}[One-dimensional errors]\label{core:record:18}
For $1024$ samples on $[-1,1)$ and Fourier cutoff $m=18$, the retained frequency indices are $-18,\ldots,18$: 37 bins including the zero-frequency bin. The measured values were
\[
\begin{array}{lrr}
\toprule
\text{method}&\RMSE&\text{maximum absolute error}\\
\midrule
\text{global Fourier truncation}&5.4023\times10^{-2}&3.4968\times10^{-1}\\
\text{known edge packet}&1.3692\times10^{-4}&1.4503\times10^{-3}\\
\text{estimated edge packet}&2.4029\times10^{-2}&5.9530\times10^{-2}\\
\bottomrule
\end{array}
\]
The simple estimator recovered the binary mask exactly in this deliberately separated two-class example. That result is a property of this test, not a universal guarantee.
\end{certificate}

\begin{figure}[htbp]
\centering
\includegraphics[width=0.84\textwidth]{06_edge_error_vs_resolution.pdf}
\caption{Error as spectral resolution increases. The known packet follows only the smooth approximation error. The estimated reconstruction reaches a floor determined by background and amplitude estimation.}
\label{fig:edge-resolution}
\end{figure}

\Needspace{12\baselineskip}
\section{Two-dimensional edge packets}\label{core:section:09}

The first image consists of a smooth background and one sharp disk. The same three logical levels must remain separate:

\begin{enumerate}[label=(\alph*)]
\item the exact theorem using the known decomposition;
\item a limited estimator designed for a positive-contrast two-class packet;
\item the open problem of recovering arbitrary edge packets under noise and model ambiguity.
\end{enumerate}

\begin{figure}[htbp]
\centering
\includegraphics[width=0.99\textwidth]{07_known_edge_packet_2d.pdf}
\caption{A sharp disk on a smooth background. Global Fourier truncation produces a bright boundary halo and oscillatory interior and exterior error. Retaining the exact edge packet leaves only the smooth-background approximation error.}
\label{fig:known-2d}
\end{figure}

\begin{certificate}[Two-dimensional errors]\label{core:record:19}
For the $192\times192$ disk experiment, a square Fourier mask retains indices $|k_x|,|k_y|\le16$, or $33^2=1089$ bins. The reported errors are
\[
\begin{array}{lrr}
\toprule
\text{method}&\text{full-image RMSE}&\text{boundary-band RMSE}\\
\midrule
\text{global Fourier truncation}&3.7643\times10^{-2}&1.8541\times10^{-1}\\
\text{known edge packet}&6.4752\times10^{-5}&1.2709\times10^{-5}\\
\text{estimated edge packet}&3.4043\times10^{-3}&8.8834\times10^{-3}\\
\bottomrule
\end{array}
\]
The mask intersection-over-union was one in this positive-contrast synthetic case. The amplitude was estimated rather than inserted from the generating model.
\end{certificate}

\begin{figure}[htbp]
\centering
\includegraphics[width=0.99\textwidth]{08_estimated_edge_packet_2d.pdf}
\caption{The limited data-driven experiment. The observed image alone is thresholded into two intensity classes, cleaned topologically, harmonically filled, spectrally smoothed, and recombined with an estimated jump.}
\label{fig:estimated-2d}
\end{figure}

\subsection{Curves, corners, and multiple amplitudes}

The exact theorem does not depend on the shape of the known edge packet. The next experiment contains a disk, a square, and a diamond with different jump amplitudes.

\begin{figure}[htbp]
\centering
\includegraphics[width=0.99\textwidth]{09_multiedge_stress_test.pdf}
\caption{Known-packet stress test with curved boundaries, corners, overlap, and different jump amplitudes. The exact packet is not a binary mask multiplied by one amplitude; it is the complete piecewise packet.}
\label{fig:multiedge}
\end{figure}

\begin{certificate}[Multiedge errors]\label{core:record:20}
The global reconstruction had RMSE $3.3740\times10^{-2}$ and maximum error $3.2807\times10^{-1}$. The known-packet reconstruction had only the remaining smooth-background error. The experiment checks the finite implementation of the theorem for a composite packet; it is not a data-driven multiregion segmentation result.
\end{certificate}

\section{Sparse wavelet packets and global Fourier modes}\label{core:section:10}

A sharp edge is spatially local but spectrally global. Wavelets often provide a more economical coordinate system once enough localized coefficients are retained. To keep the comparison explicit, the numerical study subtracts the signal mean, keeps the largest coefficients in a fixed depth-five Haar packet basis, and compares them with the same number of largest-magnitude discrete Fourier coefficients. These are coefficient-count budgets, not equal bit-storage budgets. The Fourier reconstruction uses the real part of the inverse transform; for the five reported budgets, the selected nonzero-frequency bins occur in conjugate pairs, so this operation removes only floating-point imaginary residuals.

\begin{figure}[htbp]
\centering
\includegraphics[width=0.83\textwidth]{10_sparse_wavelet_fourier.pdf}
\caption{Sparse approximation of the test signal. With very few coefficients the global Fourier basis is initially better. After the edge packets are resolved, the Haar packet error decays much faster in this example.}
\label{fig:sparse-comparison}
\end{figure}

\begin{table}[htbp]
\centering
\caption{Sparse approximation errors for the one-dimensional test signal.}
\begin{tabular}{rrr}
\toprule
\text{retained coefficients}&\text{Haar packet RMSE}&\text{Fourier RMSE}\\
\midrule
16&$1.7999\times10^{-1}$&$6.6522\times10^{-2}$\\
32&$3.8414\times10^{-2}$&$4.9153\times10^{-2}$\\
64&$1.6267\times10^{-2}$&$3.8056\times10^{-2}$\\
128&$1.7387\times10^{-3}$&$2.5631\times10^{-2}$\\
256&$1.5870\times10^{-4}$&$1.7424\times10^{-2}$\\
\bottomrule
\end{tabular}
\label{tab:sparse}
\end{table}

\begin{boundary}[Basis selection]\label{core:record:21}
This table compares two fixed finite coordinate systems for one signal. It does not prove that Haar packets are optimal, that Fourier is inferior on smooth or oscillatory data, or that a single packet tree is best for every edge geometry. The point is that the seed calculus records both systems with the same analysis, transfer, refinement, and synthesis maps.
\end{boundary}

\section{Finite adaptive packet trees}\label{core:section:11}

\begin{readerroadmap}[title={Three levels of a finite tree decision}]
First choose a complete pruning, so every terminal path is represented once and synthesis remains exact. Then assign a finite objective to those feasible choices. Additivity permits the local dynamic recursion. Finally, a positive cost gap protects the selected pruning against a declared perturbation envelope. Each step uses a different hypothesis.
\end{readerroadmap}

The preceding sparse comparison used one fixed depth-five packet basis.  A refinement tree permits a stronger choice: different branches may stop at different depths while the selected packets still form one complete coordinate system.

\begin{definition}[Complete packet pruning]\label{core:record:22}
Let $\mathcal T$ be a finite rooted refinement tree.  Every node $v$ carries a projector $P_v$, the root satisfies $P_{\varnothing}=I$, and every internal node satisfies
\[
P_v=\sum_{u\in\operatorname{ch}(v)}P_u,
\qquad
P_uP_w=0\quad(u\ne w\text{ in }\operatorname{ch}(v)).
\]
A set $\Lambda$ of nodes is a complete pruning if it is an antichain and every root-to-terminal path meets $\Lambda$ exactly once.
\end{definition}

\begin{theorem}[Exact synthesis from every complete pruning]\label{result:13}
For every complete pruning $\Lambda$,
\[
\boxed{\sum_{v\in\Lambda}P_v=I,}
\qquad
\boxed{x=\sum_{v\in\Lambda}P_vx.}
\]
Thus changing the selected tree changes the coordinates but does not introduce reconstruction error.
\end{theorem}

\begin{proof}
Starting from the root identity $P_{\varnothing}=I$, replace every unselected internal node by the exact sum of its children.  If a complete pruning contains the root it contains no other node. Otherwise its restriction to each child subtree is itself a complete pruning: every terminal path through that child must be met exactly once. Induction on the finite tree therefore justifies the substitution and its termination precisely at $\Lambda$. Child subtrees have distinct node addresses, so their selected-node sets form a disjoint union. The sum identity uses only root identity and parent-to-child conservation; sibling orthogonality gives the additional interpretation as nonoverlapping projector packets.  Applying the resulting identity to $x$ gives the synthesis formula.
\end{proof}

Let $y$ denote the observed data.  For a complete pruning $\Lambda$, consider any finite criterion
\[
J_y(\Lambda)
=
R_y(\Lambda)
+\eta L_y(\Lambda)
+\tau K(\Lambda),
\]
where $R_y$ is a representation cost, $L_y$ is a cross-packet leakage cost, and $K$ is a tree-complexity penalty.  Here $R_y,L_y,K$ take finite real values, and $\eta,\tau$ are fixed real weights (usually nonnegative in applications). No additivity is needed for the first theorem.

\begin{theorem}[Finite globally optimal stable selector]\label{result:14}
Let $\mathfrak P(\mathcal T)$ be the finite set of complete prunings of $\mathcal T$.

\begin{enumerate}[label=(\roman*)]
\item A global minimizer
\[
\Lambda_y^*\in\operatorname*{argmin}_{\Lambda\in\mathfrak P(\mathcal T)}J_y(\Lambda)
\]
always exists.
\item If the minimizer is unique and at least one competing pruning exists, its gap
\[
\Gamma_y
=
\min_{\Lambda\ne\Lambda_y^*}
\bigl(J_y(\Lambda)-J_y(\Lambda_y^*)\bigr)
\]
is attained and strictly positive. If the pruning set is a singleton, its only element minimizes every criterion; no real-valued minimum over an empty competitor set is introduced.
\item In the case with competitors, if new data $y'$ satisfy
\[
\sup_{\Lambda\in\mathfrak P(\mathcal T)}
|J_{y'}(\Lambda)-J_y(\Lambda)|
<\frac{\Gamma_y}{2},
\]
then $\Lambda_{y'}^*=\Lambda_y^*$.
\item Every selected minimizer preserves the exact synthesis identity of the preceding theorem.
\end{enumerate}
\end{theorem}

\begin{proof}
The root-only pruning exists, and every pruning is a subset of the finite node set, so the feasible set is nonempty and finite. Thus the minimum is attained. If competitors exist, their finitely many differences from a unique optimum are strictly positive, so their minimum is positive and is achieved by a competitor. In the singleton case there is no alternative selector under any objective change. In the remaining case, for any competing pruning $\Lambda\ne\Lambda_y^*$,
\[
\begin{aligned}
J_{y'}(\Lambda)-J_{y'}(\Lambda_y^*)
&\ge
J_y(\Lambda)-J_y(\Lambda_y^*)
-2\sup_{\Xi}|J_{y'}(\Xi)-J_y(\Xi)|\\
&>\Gamma_y-\Gamma_y=0.
\end{aligned}
\]
Hence the same pruning remains uniquely optimal.  Exact synthesis depends only on completeness of the pruning, not on the cost used to select it.
\end{proof}

For tree-local costs, the optimum can be found without enumerating all prunings.

\begin{theorem}[Dynamic-programming best packet basis]\label{result:15}
Suppose
\[
J_y(\Lambda)=\sum_{v\in\Lambda}c_v(y).
\]
For a node $v$, define recursively
\[
V_y(v)
=
\begin{cases}
c_v(y),&v\text{ terminal},\\[2mm]
\displaystyle
\min\left\{c_v(y),\sum_{u\in\operatorname{ch}(v)}V_y(u)\right\},
&v\text{ internal}.
\end{cases}
\]
Stop at $v$ when $c_v(y)\le\sum_uV_y(u)$, including equality; refine otherwise. This deterministic rule produces a globally optimal complete pruning. With the node costs supplied and child values cached, the evaluator uses one comparison per internal node, one visit per node, and at most $|\mathcal T|-1$ additions. Its arithmetic-and-visit count is therefore at most $3|\mathcal T|$. This is a unit-cost statement about the specified evaluator; it excludes evaluation of the local costs, bit complexity of scalar arithmetic, and materialization or sorting of packet-address strings.
\end{theorem}

\begin{proof}
Proceed upward from the terminal nodes.  Every complete pruning below an internal node $v$ has exactly two possible forms: it either stops at $v$, with cost $c_v(y)$, or it refines and is the union of complete child prunings, whose minimum total cost is the sum of the already computed child minima.  The recursion therefore gives the exact minimum on every subtree, in particular at the root. This characterization also proves both directions of the correspondence between recursive cuts and complete covering antichains, for arbitrary positive finite branching, including unary branches. Summing cached child values charges at most one addition per child edge; a finite rooted tree has $|\mathcal T|-1$ such edges. The number of internal nodes is at most $|\mathcal T|$, which gives the stated work bound. Ties do not affect optimality, but the stop-on-ties rule makes the selected pruning unambiguous.
\end{proof}

\begin{corollary}[Local perturbation certificate]\label{result:16}
Assume
\[
|c_v(y')-c_v(y)|\le\varepsilon_v.
\]
Define
\[
E(v)
=
\begin{cases}
\varepsilon_v,&v\text{ terminal},\\[1mm]
\displaystyle
\max\left\{\varepsilon_v,
\sum_{u\in\operatorname{ch}(v)}E(u)\right\},&v\text{ internal}.
\end{cases}
\]
For the complete prunings $\mathfrak P_v$ below $v$, the envelope is exactly
\[
E(v)=\max_{\Lambda\in\mathfrak P_v}\sum_{a\in\Lambda}\varepsilon_a.
\]
The maximum is attained. Every complete-pruning cost changes by at most $E(\varnothing)$, and the two globally optimal values differ by at most this quantity. Consequently, for a unique optimum with competitors,
\[
\boxed{\Gamma_y>2E(\varnothing)}
\]
is a deterministic certificate that the globally optimal pruning is unchanged.  A singleton pruning set is stable without a gap condition. Assume the random node costs and the joint local-error event are measurable. If a noise model supplies \emph{all} local bounds simultaneously on one event of probability at least $1-\delta$, and the displayed strict-gap condition holds, the same selector is recovered with probability at least $1-\delta$. Individual nodewise probabilities alone do not give this joint confidence without an additional argument, such as a union bound.
\end{corollary}

\begin{proof}
For any pruning of the subtree rooted at $v$, either $v$ itself is selected, producing perturbation at most $\varepsilon_v$, or the pruning is the disjoint union of child prunings, producing perturbation at most the sum of the child envelopes.  The same decomposition applied to the nonnegative weights $\varepsilon_a$ proves equality with the maximum weight: at each node choose either the stop option or a maximizing pruning in each child, whichever gives the larger total. Finiteness ensures these choices attain their maxima. Evaluating each cost at a minimizer of the other cost bounds the difference of the two optimal values by the uniform envelope. The selector statement follows from the strict-gap theorem. The finite deterministic comparisons defining the selector make its preservation event measurable when the random node costs are measurable. The joint local-error event is contained in the selector-preservation event, so monotonicity of probability proves the probabilistic assertion.
\end{proof}

\begin{figure}[htbp]
\centering
\includegraphics[width=0.97\textwidth]{17_adaptive_packet_tree.pdf}
\caption{A globally optimal mixed-depth Haar packet pruning.  The selected antichain covers the full tree, so its synthesis is exact even though different branches stop at different depths.}
\label{fig:adaptive-tree}
\end{figure}

\begin{figure}[htbp]
\centering
\includegraphics[width=0.97\textwidth]{18_adaptive_tree_stability.pdf}
\caption{The dynamic-programming optimum agrees with exhaustive enumeration.  A positive gap larger than twice the declared perturbation envelope certifies that the selected tree is stable.}
\label{fig:adaptive-stability}
\end{figure}

\begin{certificate}[Adaptive-tree optimum]\label{core:record:27}
For the depth-five example, the complete binary tree has $458{,}330$ admissible prunings.  Dynamic programming and exhaustive enumeration returned the same minimum with residual $0$.  The selected pruning was
\[
\{000,001,0100,0101,011,1\},
\]
and mixed-depth synthesis reconstructed the signal with maximum residual $8.88\times10^{-16}$.  The optimum gap was $3.1909\times10^{-4}$, while twice the declared perturbation envelope was $6.4\times10^{-5}$, so the finite selector is certified stable.
\end{certificate}

\begin{boundary}[Finite versus asymptotic selection]\label{core:record:28}
The theorem resolves the finite adaptive-tree existence, exact-synthesis, global-optimality, and margin-stability problem.  It does not by itself prove minimax rates, statistical consistency on an infinite tree, or a universal choice of representation and leakage costs for every signal class.
\end{boundary}


\section{Nonlinear signal packets}\label{core:section:12}

Linear time invariance is not required for exact packet coordinates. Let $A$ and $B$ be formal power series over a commutative coefficient ring (in particular $\bbC$). Set
\[
A(z)=\sum_{n\ge0}a_nz^n,
\qquad
A_i=P_{N,i}A.
\]
All phase indices range from $0$ to $N-1$. For another representation $B(z)$, multiplication gives a nonlinear or bilinear interaction.

\begin{theorem}[Cyclic product law]\label{result:17}
\[
\boxed{
P_{N,k}(AB)
=
\sum_{i+j\equiv k\, (\mathrm{mod}\,N)}A_iB_j.
}
\]
For $q\in\bbN$, the $q$th-power identity is
\[
P_{N,k}(A^q)
=
\sum_{i_1+\cdots+i_q\equiv k}
A_{i_1}\cdots A_{i_q}.
\]
\end{theorem}

\begin{proof}
Every coefficient in $A_iB_j$ has degree congruent to $i+j$ modulo $N$. At any fixed degree $n$, the Cauchy coefficient sum is finite. Sorting these terms by the residues of their two indices gives the first identity coefficient by coefficient. For the power identity, use induction on $q$ and the finite product law. At $q=0$, the index set contains the single empty tuple, its product is $1$, and its residue sum is $0$; hence the right side is $1$ for $k=0$ and $0$ otherwise, as required. No convergence is needed for these formal identities. For analytic power series, the same identities hold wherever the series are absolutely convergent and their Cauchy products may be evaluated.
\end{proof}

\begin{figure}[htbp]
\centering
\includegraphics[width=0.76\textwidth]{14_nonlinear_phase_coupling.pdf}
\caption{Norms of pairwise product interactions for five residue phases. The label in each cell gives the output phase $(i+j)\bmod5$.}
\label{fig:nonlinear-coupling}
\end{figure}

For a finite nonlinear shift operator
\[
F(x)[n]=x[n]+\alpha x[n]x[n-1],
\]
the exact packet coordinates are simply
\[
F_r((x_s)_s)=P_{N,r}F\!\left(\sum_sP_{N,s}x\right).
\]
Analysis followed by this packet system and synthesis reproduces the direct nonlinear output exactly. The numerical residual in the companion implementation was zero to machine representation.

\section{Smith coordinates as cyclic signal coordinates}\label{core:section:13}

Let $Z_0>0$ be a reference impedance and
\[
z=\frac{Z}{Z_0}=r+\ii x.
\]
The reflection coefficient is
\[
\boxed{
\Gam=\frac{z-1}{z+1},
\qquad
z=\frac{1+\Gam}{1-\Gam}.
}
\]

\begin{classical}[Cayley geometry]\label{core:record:30}
The map sends the finite passive half-plane $r\ge0$ onto the closed unit disk with $\Gam=1$ omitted. Adjoining the point at infinity in impedance adds that missing boundary point; the open right half-plane maps conformally onto the open disk. Constant resistance and reactance become the familiar Smith circles. This conformal geometry is classical and is not replaced by polyphase refinement.
\end{classical}

Writing $\Gam=u+\ii v$, the grid equations are
\[
\left(u-\frac{r}{1+r}\right)^2+v^2=\frac1{(1+r)^2},
\]
and, for $x\ne0$,
\[
(u-1)^2+\left(v-\frac1x\right)^2=\frac1{x^2}.
\]

Fix
\[
\omega_N=e^{2\pi\ii/N}
\]
and rotate the reflection coefficient by
\[
R_{N,k}(\Gam)=\omega_N^k\Gam.
\]

\begin{theorem}[Polyphase M\"obius action]\label{result:18}
Regard normalized impedance as the projective line $\mathbb P^1(\bbC)=\bbC\cup\{\infty\}$. For $N\ge1$, the pullback is a globally defined projective action. In the finite chart, wherever the denominator is nonzero, it is
\[
\boxed{
T_{N,k}(z)
=
\frac{(1+\omega_N^k)z+(1-\omega_N^k)}
{(1-\omega_N^k)z+(1+\omega_N^k)}.
}
\]
As maps on $\mathbb P^1(\bbC)$ it satisfies
\[
T_{N,k}\circ T_{N,\ell}=T_{N,k+\ell\, (\mathrm{mod}\,N)}.
\]
\end{theorem}

\begin{proof}
Let $C(z)=(z-1)/(z+1)$. Then
\[
T_{N,k}=C^{-1}\circ R_{N,k}\circ C.
\]
For a division-free construction, represent a point by a nonzero pair $[u:v]$, identifying $[cu:cv]=[u:v]$ for $c\ne0$. With $w=\omega_N^k$, use the matrix
\[
A(w)=\frac12\begin{pmatrix}1+w&1-w\\1-w&1+w\end{pmatrix}.
\]
It has determinant $w\ne0$, preserves nonzero pairs, and commutes with rescaling, so it induces a projective map. Direct multiplication gives $A(w)A(w')=A(ww')$ and $A(1)=I$. Since $\omega_N^N=1$, this proves the cyclic composition law modulo $N$. The Cayley coordinate of $[u:v]$ is $[u-v:u+v]$, and applying $A(w)$ changes this pair to $[w(u-v):u+v]$. This proves the conjugacy also at poles. For $[z:1]$, taking the coordinate ratio gives the displayed finite-chart formula.
\end{proof}

\begin{corollary}[Impedance--admittance is the two-phase split]\label{result:19}
For $N=2$ and finite $z\ne0$,
\[
\boxed{T_{2,1}(z)=z^{-1}.}
\]
On the full projective line it exchanges $0$ and $\infty$. Thus a half-turn of the reflection disk exchanges normalized impedance and admittance.
\end{corollary}
\begin{proof}
$\omega_2=\exp(\pi\ii)=-1$, so $A(\omega_2)[u:v]=[v:u]$. The finite ratio is $1/z$ when $z\ne0$; the coordinate swap also gives $[0:1]\leftrightarrow[1:0]$. No ordinary division by zero is used to represent infinity.
\end{proof}

\begin{figure}[htbp]
\centering
\includegraphics[width=0.88\textwidth]{11_polyphase_smith_chart.pdf}
\caption{The classical Smith grid equipped with sixteen exact phase rays and four radial seed orbits. The Cayley transform creates the chart; the projector calculus supplies the nested cyclic coordinate system.}
\label{fig:smith-chart}
\end{figure}

\subsection{Radial seeds and physical phase}

For $0\le\rho<1$ and angle $\theta$,
\[
\Gam=\rho e^{\ii\theta}
\]
gives
\[
z(\rho,\theta)
=\frac{1+\rho e^{\ii\theta}}{1-\rho e^{\ii\theta}},
\]
with
\[
r(\rho,\theta)=\frac{1-\rho^2}{1-2\rho\cos\theta+\rho^2},
\qquad
x(\rho,\theta)=\frac{2\rho\sin\theta}{1-2\rho\cos\theta+\rho^2}.
\]
A fixed $\rho$ is a constant-mismatch circle. Refining $N$ increases angular resolution without changing $\rho$.

\begin{figure}[htbp]
\centering
\includegraphics[width=0.78\textwidth]{12_smith_nested_refinement.pdf}
\caption{The same reflection circle at phase orders $2,4,8,16,32$. Each coarse phase ray appears in every finer grid whose order is a multiple of its order. The exact residue-projector refinement law concerns coefficient packets, not sums of impedance points.}
\label{fig:smith-refinement}
\end{figure}

For a lossless transmission line,
\[
\Gam_{\rm in}=\Gam_L e^{-2\ii\beta\ell}.
\]
One $N$-phase step corresponds to a line-length increment $\lambda/(2N)$, up to the chosen orientation convention. Hence the phase address has a direct electrical interpretation.

\section{Dyadic--odd phase factorization and boundary weighting}\label{core:section:14}

Write
\[
N=2^sM,
\qquad M\text{ odd}.
\]
The cyclic index splits by the Chinese remainder theorem:
\[
\bbZ/N\bbZ
\cong
\bbZ/2^s\bbZ\times\bbZ/M\bbZ.
\]

\begin{proposition}[Dyadic--odd packet factorization]\label{result:20}
After identifying the residue labels by the Chinese remainder map, the phase projector factors as
\[
\boxed{
P_{N,j}=P_{2^s,j_2}P_{M,j_o}.
}
\]
\end{proposition}

\begin{proof}
The two congruence conditions are independent because $2^s$ and $M$ are coprime. Their indicator functions multiply to the indicator of the unique residue class modulo $N$ satisfying both.
\end{proof}

The first dyadic step is exceptional in Smith coordinates because it is the inversion $z\mapsto z^{-1}$. Higher dyadic levels refine the open/short duality into quarter-, eighth-, and finer phase sectors.

Away from its pole $\Gam=1$, the inverse Cayley map has derivative
\[
\frac{dz}{d\Gam}=\frac{2}{(1-\Gam)^2},
\]
while normalized admittance has derivative proportional to $(1+\Gam)^{-2}$. Away from $\Gam=\pm1$, this motivates the symmetric sensitivity
\[
\boxed{
w_{\rm Smith}(\Gam)
=
\frac1{|1-\Gam|^2}
+
\frac1{|1+\Gam|^2}.
}
\]

\begin{figure}[htbp]
\centering
\includegraphics[width=0.78\textwidth]{13_smith_sensitivity.pdf}
\caption{The proposed symmetric refinement weight becomes large near the open- and short-circuit boundary points. It is a condition-based adaptive coordinate, not an asserted optimal sampling law.}
\label{fig:smith-weight}
\end{figure}

\begin{boundary}[Weighted refinement]\label{core:record:34}
The sensitivity formula follows directly from the exact coordinate derivatives. Using it to choose an optimal adaptive grid, interpolation rule, or measurement design is a separate numerical problem. The release visualizes the weight but does not prove optimality.
\end{boundary}

\section{Boundary-localized error and claim control}\label{core:section:15}

The strongest visual distinction between global Fourier reconstruction and exact packet retention is seen in a band around the discontinuity.

\begin{figure}[htbp]
\centering
\includegraphics[width=0.94\textwidth]{15_boundary_leakage.pdf}
\caption{Error restricted to a narrow boundary band. Global truncation concentrates large error on the jump. The known-packet reconstruction has only the small smooth-packet error in the same band.}
\label{fig:boundary-error}
\end{figure}

\begin{figure}[htbp]
\centering
\includegraphics[width=0.90\textwidth]{16_claim_status.pdf}
\caption{Proof and claim boundary. Finite adaptive-tree selection is now proved; the remaining open rows concern universal recovery and infinite-tree or asymptotic selection.}
\label{fig:claim-status}
\end{figure}

\begin{conjecture}[Edge-packet approximation class]\label{core:record:35}
There should exist natural function spaces in which a signal is represented as
\[
f=s+\sum_q a_q\mathbf1_{\Omega_q},
\]
with separate norms for smooth approximation, boundary geometry, packet amplitude, and packet identification. A projector-resolved approximation theory on such spaces should yield rates that do not mix smooth truncation error with discontinuity-location error.
\end{conjecture}


\section{Implementation and reproducibility}\label{core:section:16}

The companion package is organized by mathematical layer:

\begin{longtable}{L{0.27\textwidth}L{0.64\textwidth}}
\caption{Numerical companion modules.}\label{tab:modules}\\
\toprule
\textbf{module}&\textbf{implemented content}\\
\midrule
\endfirsthead
\toprule
\textbf{module}&\textbf{implemented content}\\
\midrule
\endhead
\texttt{polyphase.py}&residue projectors, analysis/synthesis, refinement, exact FIR transfer kernels, polyphase convolution\\
\texttt{filterbanks.py}&Haar analysis/synthesis, complete Haar packet trees, energy accounting, predict--update lifting\\
\texttt{adaptive.py}&complete-pruning enumeration, exact mixed-depth synthesis, dynamic-programming optimization, and stability certificates\\
\texttt{edges.py}&one- and two-dimensional Fourier truncation, known-packet reconstruction, narrow two-class estimators, harmonic filling, error metrics\\
\texttt{smith.py}&Cayley coordinates, cyclic M\"obius action, Smith circles, sensitivity weight, coefficient projectors\\
\texttt{nonlinear.py}&coefficient packets, cyclic product convolution, nonlinear shifted packet reconstruction\\
\texttt{experiments.py}&all machine-readable numerical certificates and tabular summaries\\
\texttt{visuals.py}&the eighteen figure pairs used in the paper\\
\bottomrule
\end{longtable}

The September 2026 release records recovery and rerunning of the original companion source: its 33 tests passed, and all 18 figure pairs and numerical tables were regenerated. The Lean package contributes 20 independent exact-model tests and 16 tests of the verifier itself; these account for another 36 passing tests. The exact-model suites count 305,933 assertions. These counts are separate from the 225 Lean theorem targets.

All numerical computations are CPU-based. Floating-point residuals measure finite implementations and may vary at rounding level with the scientific libraries. The universal mathematical statements are supported by the Lean proofs; the Python modules are not compiled or verified by Lean. The adaptive-tree certificate uses floating-point costs and an exhaustive finite comparison, so its displayed gap is numerical evidence rather than an interval-arithmetic proof about those computed costs. Appendix~\ref{sec:formalization} explains the verification boundary.

\begin{certificate}[Finite numerical residuals]\label{core:record:36}
The largest recorded residuals were
\[
\begin{array}{lr}
\toprule
\text{identity}&\text{residual}\\
\midrule
\text{polyphase reconstruction}&0\\
\text{nested refinement}&0\\
\text{polyphase convolution}&4.44\times10^{-16}\\
\text{Haar reconstruction}&8.88\times10^{-16}\\
\text{wavelet-packet reconstruction}&2.22\times10^{-15}\\
\text{adaptive-tree optimum versus exhaustive search}&0\\
\text{adaptive mixed-depth synthesis}&8.88\times10^{-16}\\
\text{lifting reconstruction}&4.44\times10^{-16}\\
\text{Smith rotation conjugacy}&3.34\times10^{-16}\\
\text{Smith impedance--admittance duality}&0\\
\text{nonlinear product-packet reconstruction}&1.33\times10^{-15}\\
\text{nonlinear shifted-filter reconstruction}&0\\
\bottomrule
\end{array}
\]
The Haar and wavelet energy residuals were $5.68\times10^{-14}$ and $2.27\times10^{-13}$, respectively.
\end{certificate}

\section{Conclusions}\label{core:section:17}

The signal realization makes the seed, address, packet, transfer, refinement, and synthesis operations concrete. The objects and laws can be read directly:
\[
\boxed{
\begin{aligned}
\text{seed}&=\delta_0,\\
\text{address}&=S^n,\\
\text{packet}&=P_{N,r}x,\\
\text{transfer}&=P_{N,k}HP_{N,s},\\
\text{refinement}&=P_{N,r}=\sum_{s\succ r}P_{M,s},\\
\text{synthesis}&=\sum_rP_{N,r}x,\\
\text{nonlinearity}&=P_{N,r}F\left(\sum_sP_{N,s}x\right).
\end{aligned}
}
\]
Filter banks, wavelet packets, adaptive complete prunings, edge packets, and Smith coordinates inhabit the same operator grammar, but their claims differ. Polyphase and filter-bank identities are exact. Wavelet perfect reconstruction is exact. The known-edge theorem is exact conditional on the decomposition. Estimated-edge performance depends on a detection model. The Smith chart remains the classical Cayley geometry equipped with an exact cyclic refinement algebra. Keeping those distinctions explicit is part of the calculus, not an editorial afterthought.

The next paper in the series turns from discrete signal transfer to semigroups, resolvents, conditioning, stochastic evolution, fractional powers, and PDE kernels.

\appendix
\section{Algorithm statements}\label{core:section:18}

\subsection{Exact polyphase filtering}

Given $x$, $h$, and $N$:

\begin{enumerate}
\item Form $x_s[n]=x[Nn+s]$ for $0\le s<N$.
\item Form every transfer kernel $h_{ks}[m]=h[Nm+k-s]$.
\item Compute
\[
y_k[n]=\sum_s\sum_mh_{ks}[m]x_s[n-m].
\]
\item Interleave the $y_k$ channels.
\end{enumerate}

The algorithm is algebraically identical to direct convolution; its purpose is to expose independent and interacting packets for block computation and refinement.

\subsection{Conditional known-edge reconstruction}

Given a decomposition $f=s+j$:

\begin{enumerate}
\item Apply the selected smooth approximation $Q_m$ to $s$ only.
\item Retain $j$ exactly.
\item Return $Q_ms+j$.
\end{enumerate}

The error certificate is $\|s-Q_ms\|$.

\subsection{Limited positive-contrast estimator}

The included finite estimator does the following:

\begin{enumerate}
\item Split the observed values into two classes by an Otsu threshold.
\item Keep coherent connected components and fill holes.
\item Harmonically fill the selected packet to estimate the missing background.
\item Smooth the filled background spectrally.
\item Estimate one packet amplitude by a robust median residual.
\item Recombine the estimated smooth and edge packets.
\end{enumerate}

This procedure is useful for the supplied separated synthetic case. It assumes a positive-contrast two-class packet and is not promoted to a general edge theorem.

\section{Numerical record}\label{core:section:19}

\begin{table}[H]
\centering
\caption{Main edge-study measurements.}
\begin{tabular}{lrrr}
\toprule
\text{case}&\text{global RMSE}&\text{known-packet RMSE}&\text{estimated RMSE}\\
\midrule
1D rectangular packet&$5.4023\times10^{-2}$&$1.3692\times10^{-4}$&$2.4029\times10^{-2}$\\
2D disk packet&$3.7643\times10^{-2}$&$6.4752\times10^{-5}$&$3.4043\times10^{-3}$\\
2D multiedge packet&$3.3740\times10^{-2}$&\multicolumn{2}{c}{known composite packet only}\\
\bottomrule
\end{tabular}
\end{table}

\begin{table}[H]
\centering
\caption{Main software artifacts.}
\begin{tabular}{ll}
\toprule
\text{artifact}&\text{purpose}\\
\midrule
\texttt{results/numerical\_results.json}&complete machine-readable certificate record\\
\texttt{tables/numerical\_summary.csv}&flat summary of errors and residuals\\
\texttt{tables/sparse\_representation.csv}&wavelet/Fourier coefficient study\\
\texttt{tests/}&automated finite identity and estimator tests\\
\texttt{figures/}&PNG and vector-PDF figure pairs\\
\texttt{VERIFY\_PAPER.py}&combined manuscript, numerical, and Lean checks\\
\bottomrule
\end{tabular}
\end{table}

\section{Notation crosswalk}\label{core:section:20}

\begin{longtable}{L{0.20\textwidth}L{0.33\textwidth}L{0.37\textwidth}}
\caption{Notation crosswalk.}\label{tab:notation}\\
\toprule
\textbf{common symbol}&\textbf{signal realization}&\textbf{interpretation}\\
\midrule
\endfirsthead
\toprule
\textbf{common symbol}&\textbf{signal realization}&\textbf{interpretation}\\
\midrule
\endhead
$s$&$\delta_0$&primitive signal seed\\
$A_n$&$S^n$&address operator\\
$P_{N,r}$&residue mask&phase packet projector\\
$T_{ks}$&$P_{N,k}HP_{N,s}$&input-to-output phase transfer\\
$E,R$&analysis and synthesis matrices&coordinate conversion and reconstruction\\
$W(\gamma)$&filter or path weight&refinement or implementation cost\\
$\Delta(F)$&coarse/fine signal defect&aliasing, leakage, or estimator inconsistency\\
$\operatorname{Hol}(\gamma)$&closed phase transfer&cycle consistency in a network\\
$\Gam$&reflection coefficient&Smith-disk coordinate\\
\bottomrule
\end{longtable}

\clearpage
\section{Formalization and reproducibility record}\label{core:section:21}\label{sec:formalization}

\subsection{What has been checked}

The preserved Part I manuscript has twenty theorem-like statements: Theorems, Propositions, and Corollaries, counted together. The Lean project contains 225 explicit supporting theorem targets in 23 proof modules. All twenty statements have an explicit scope row and principal declaration targets. The September 2026 release records a successful build and axiom audit with warnings treated as errors. No fresh Lean build is asserted for this October final edition. The audit inspected all 827 declarations in the project namespace, including definitions and automatically generated declarations. The latter is an implementation census, not a count of 827 mathematical theorems.

The accepted axiom dependencies are \texttt{propext}, \texttt{Classical.choice}, and \texttt{Quot.sound}. No admitted proof or additional project axiom is permitted. The verifier checks the imported modules, registered theorem inventory, dependency pins, source hashes, and coverage of the twenty-row scope contract. Its successful status is \texttt{LEAN\_CHECKS\_PASS}. The correspondence between a mathematical sentence in this PDF and its Lean type is a manual mathematical review, not an automated semantic-equivalence test.

The source-level formalization was completed before this manuscript revision. The revision integrates those checked hypotheses and conclusions into the paper itself; it does not increase the Lean theorem count by counting editorial changes as new proofs. Eleven foundational formal-power-series proofs in \texttt{CyclicSeries} are adaptations of Paper IV and are identified as such in the provenance record.

\subsection{Result-by-result correspondence}

The following table summarizes the reviewed statement boundary. Full qualified declaration names, source locations, and supplementary targets appear in \texttt{MANUSCRIPT\_TARGET\_MAP.md} and the machine-readable scope contract. Every row below has supporting kernel-checked declarations; numerical certificates, classical-input paragraphs, remarks, and conjectures are not additional theorem rows.

\begingroup\small
\begin{longtable}{L{0.09\textwidth}L{0.28\textwidth}L{0.53\textwidth}}
\caption{The twenty manuscript results and their Lean statement scopes.}\label{tab:lean-map}\\
\toprule
\textbf{Result}&\textbf{Principal modules}&\textbf{Checked mathematical content}\\
\midrule\endfirsthead
\toprule
\textbf{Result}&\textbf{Principal modules}&\textbf{Checked mathematical content}\\
\midrule\endhead
\ref{result:01}&HilbertPackets&Complex $\ell^2$ impulse expansion, finite-support identity, unitary shifts, and closed shift-generated span.\\
\ref{result:02}&HilbertPackets&Bounded self-adjoint residue projectors, pairwise annihilation, completeness, and orthogonal ranges.\\
\ref{result:03}&IntegerPolyphase; HilbertPackets&Signed quotient/remainder coordinates, both inverse identities, and the joint Hilbert norm of the channels.\\
\ref{result:04}&HilbertPackets&Literal bounded-projector refinement for positive divisible moduli.\\
\ref{result:05}&FIRTransfer; ZTransform&Finite-support tap reindexing for all bilateral inputs; numerical channel transforms under weighted absolute summability.\\
\ref{result:06}&FilterBank; FiniteTransforms&Finite Laurent/FIR bank composition, faithful action on sequences, gain-delay reconstruction criterion, and the $Nd$ original-sample delay.\\
\ref{result:07}&HaarChannels; HaarHilbert&Literal even/odd input samples, normalized complex low/detail channels, and $\ell^2$ energy equality.\\
\ref{result:08}&Lifting; FiniteTransforms&Predict/update inverse maps, determinant one, and both flat finite-signal inverse identities.\\
\ref{result:09}&WaveletPackets&Every finite depth of literal Haar packet operators, energy equality, and reconstruction using actual Hilbert adjoints.\\
\ref{result:10}&KnownEdge&Exact cancellation of one retained packet, and equality of approximation errors.\\
\ref{result:11}&KnownEdge&Finite families of retained packets, with no geometry assumption on the packets.\\
\ref{result:12}&MaskPerturbation&Measurable-mask symmetric-difference norm, extended finite-$p$ error bound, pointwise and essential supremum bounds, and smooth-background error.\\
\ref{result:13}&CompletePrunings; PruningSums&Recursive cuts correspond to complete covering antichains; literal finite node sums preserve synthesis.\\
\ref{result:14}&PruningSums&Finite arbitrary-objective minimization, attained positive gap with competitors, strict stability, and the singleton case.\\
\ref{result:15}&FiniteAdaptive; PruningAlgorithm&Deterministic stop-on-ties selector, global optimum, arbitrary positive finite branching, and the stated unit-operation count.\\
\ref{result:16}&PruningSums; PruningAlgorithm&Attained maximum error envelope, optimal-value perturbation, strict-gap selector stability, and probability transfer from a joint event.\\
\ref{result:17}&CyclicSeries&Formal Cauchy-product residue convolution and the finite-tuple power identity, including the empty product at exponent zero.\\
\ref{result:18}&SmithAction; SmithCyclic&Nonzero homogeneous pairs modulo scaling, Cayley conjugacy, the actual exponential phase, and cyclic composition on the projective line.\\
\ref{result:19}&SmithCyclic&Two-phase coordinate swap, $0\leftrightarrow\infty$, and finite impedance inversion when $z\ne0$.\\
\ref{result:20}&CRT&Signed integer residue-mask factorization for coprime moduli, including the dyadic--odd case.\\
\bottomrule
\end{longtable}
\endgroup

The remaining foundational modules \texttt{Packets}, \texttt{FiniteSelection}, and \texttt{BinaryAdaptive} provide generic coordinate, finite-minimum, and binary-tree lemmas used or retained alongside the concrete constructions. The Hilbert target is the square-summable space on a product of indices; it is not the supremum-norm function space that Lean would otherwise assign to an unqualified finite family of channels.

\subsection{Three bridges that matter for the statement}

\paragraph{Finite filters and numerical transforms.}
The filter-bank proof is an algebraic theorem about finite Laurent-polynomial kernels acting faithfully on bilateral sequences. Equality at a single complex value of $z$ would be insufficient to prove a polynomial identity; the impulse argument avoids that error. The separate transform theorem proves the finite-filter action agrees with evaluation wherever the weighted input series converge absolutely. An arbitrary bilateral $\ell^2$ sequence need not have an absolutely convergent numerical $z$-transform. The all-sequence time-domain theorem makes no such assumption.

\paragraph{Recursive trees and set-valued prunings.}
The formal tree permits any positive finite arity at an internal node, including arity one. A node has a position in the tree, so equal data labels on distinct nodes do not identify the nodes. A recursive cut is either a stop or a collection of child cuts. Its equivalence with a covering antichain is proved in both directions. The cost, envelope, and operator sums are then shown to equal actual finite sums over the selected node set. This connects the recursive optimizer to the manuscript's set-valued definition instead of assuming that connection.

\paragraph{Hilbert operators and finite experiments.}
The coordinate maps, literal Haar channels, and packet operators are constructed on complex $\ell^2$. Square summability and boundedness are established before adjoints are used. The finite Python experiments separately check indexing, reconstruction, and energy at finite lengths. They illustrate these operators but are not the objects executed by the Lean kernel. Likewise, the Python Smith routine evaluates only finite charts; the Lean projective construction supplies the pole-safe universal statement.

\subsection{September 2026 checks and numerical conventions}

\begin{center}
\begin{tabular}{lr}
\toprule
\textbf{Check family}&\textbf{Accepted count}\\
\midrule
Original numerical companion tests&33\\
Independent exact-model tests&20\\
Verifier positive and negative controls&16\\
Total tests in these three families&69\\
Exact assertions inside the independent model suites&305,933\\
Registered Lean theorem targets&225\\
Lean proof modules&23\\
Regenerated figure pairs&18\\
\bottomrule
\end{tabular}
\end{center}

The September 2026 record states that all scalar entries in the recovered numerical-results JSON reproduced exactly in its environment. This includes the 458,330-pruning depth-five comparison, its selected six-node pruning, the optimum gap, and the reconstruction residuals. The five Fourier sparse-approximation budgets were also checked for conjugate-pair closure. The revised figures and tables use the regenerated data. This exact replay of the stored floating-point numbers is stronger than a visual similarity check, but it is not a promise of bitwise equality on every platform.

For the one-dimensional low-pass experiment, ``cutoff 18'' means 37 signed frequency bins, not 18 retained coefficients. The disk experiment uses a square $33\times33$ frequency mask, not a radial cutoff. The sparse comparison uses largest-coefficient budgets instead of either low-pass convention. The full-image and boundary-band RMSEs are normalized by the number of samples in the respective comparison set. The supplied inpainting routines perform a fixed finite number of periodic neighbor-averaging iterations; ``harmonic filling'' describes that iteration, not a separately proved exact continuum Dirichlet solver.

The independent exact-model suite checks signed polyphase/FIR addresses, exact lifting and normalized Haar arithmetic, mask bounds, and exhaustive finite-tree objectives. Its tree census includes all 626 ordered rooted tree shapes with one through eight nodes and 129,958 node subsets. These finite checks can catch implementation or convention errors; they are not substitutes for the universal Lean proofs. The verifier controls exercise failure paths, including incomplete scope contracts and unregistered proof targets. Mandatory verification conditions use explicit exceptions rather than Python assertions, so optimized Python cannot silently remove them.

\subsection{Pinned environment and rerunning the release}

The Lean toolchain is \texttt{leanprover/lean4:v4.34.0-rc2}. The project pins mathlib to
\begin{center}\small\texttt{aeb72f1e9981b6d35b426b16dca7759be3e7ba06}.\end{center}
The Lake manifest records the remaining dependency revisions. The numerical environment versions are recorded in \texttt{verification/numerical\_environment.json}. Source and artifact hashes appear in \texttt{SHA256SUMS.txt}. The original manuscript, source, numerical data, and figures remain available under provenance for comparison; the expanded PDF and its TeX are the manuscript of that September 2026 release.

From the preserved release root (\texttt{source/complete/original\_snapshot} in the final package), the combined verification entry point is:
\begin{lstlisting}[language=bash,numbers=none]
python3 VERIFY_PAPER.py
\end{lstlisting}
This runs the Lean verifier, companion tests, numerical regeneration/comparison, and manuscript consistency checks. Lean, Lake, the pinned dependencies, the Python requirements, and a TeX installation must be available for the corresponding full checks. Missing prerequisites or failed checks return failure; static checks alone never report kernel verification. The individual numerical generation command is
\begin{lstlisting}[language=bash,numbers=none]
cd companion
python3 run_all.py
python3 -m pytest -q tests
\end{lstlisting}
The TeX source resolves the eighteen vector figures from the manuscript's figure directory. The checked build uses repeated pdfLaTeX passes, then text, cross-reference, and page-render inspection. Different pdfTeX versions can produce different binary PDF hashes even when the rendered document is unchanged.

In the September 2026 release environment, a disclosed runtime adapter was needed for Lean executable-path discovery: it redirects a read of the current process's numeric \texttt{/proc} executable link to \texttt{/proc/self/exe}. It does not modify the Lean binary, elaborator, kernel, axioms, or proof terms. The adapter source and invocation are retained with the verification evidence; ordinary installations can run the pinned toolchain directly.

\subsection{What remains open}

The September 2026 finite theorem scope is complete in the stated sense: all twenty theorem-like statements have reviewed, accepted formal counterparts, and that release records a rerun of the numerical companion. The two edge-recovery conjectures remain open research directions. No universal segmentation theorem, infinite-tree consistency result, minimax rate, or optimal Smith sampling theorem is inferred from the finite results. Classical Smith-grid and transmission-line formulae provide application context; those physical and historical assertions are not additional Lean-certified theorem rows.

Across the available Papers I--V manuscripts there are 168 theorem-like results. Their release maps contain 1,456 raw theorem targets; removing two shared qualified names and eleven known Paper-IV-to-V adaptations leaves 1,443 after those known duplicate removals. This is a source census, not a claim of that many independent or novel mathematical results. Papers I--IV were counted from their release material and bundled audit records; only Paper V was freshly built for that September 2026 release. These historical cross-paper counts are not a current census of the October final companions.

\clearpage
\part*{Part II: Extended mathematical development}\addcontentsline{toc}{part}{Part II: Extended mathematical development}
\begin{center}\large Multiresolution Signal Calculus, Singular Packets, and Stable Reconstruction\end{center}
\noindent This part preserves the complete standalone mathematical branch, with the explicit audit repairs recorded in the accompanying source ledger. Its sixty-three proofs are prose proofs; no additional Lean coverage is asserted. Section and result numbers carry the prefix E to distinguish this scope from Part I.
\setcounter{section}{0}\setcounter{figure}{0}\setcounter{table}{0}
\renewcommand{\thesection}{E\arabic{section}}
\renewcommand{\theHsection}{extension.\arabic{section}}
\renewcommand{\thefigure}{E\arabic{figure}}
\renewcommand{\theHfigure}{extension.\arabic{figure}}
\renewcommand{\thetable}{E\arabic{table}}
\renewcommand{\theHtable}{extension.\arabic{table}}
\graphicspath{{extension_figures/}}

\begin{abstract}
Signals on a lattice admit exact coordinates indexed by the residue classes of every finite index sublattice. Convolution becomes a matrix of packet transfers, while filter banks and packet trees become exact changes of coordinates. We develop this structure on sequence spaces, prove quantitative stability for imperfect and perturbed banks, and give an exact dynamic recursion for adaptive packet trees. A singular signal is written as a smooth component plus step or truncated power packets that record jumps in the signal and its derivatives. Retaining the correct packets leaves only the smooth approximation error. Estimated packets are governed by exact formulas involving value error, region overlap, boundary displacement, and the norm of the smooth approximation operator. The main recovery result treats separated jumps sampled on a uniform grid. An explicit contrast gap gives exact recovery of every jump cell, grid scale location accuracy, amplitude control, and finite norm consistency. Probability bounds cover sub Gaussian noise and repeated measurements. Haar coefficients, two shifted grids, and annihilating differences give complementary localization results. The final sections treat Gibbs oscillation, nonlinear residue interactions, numerical illustrations, and cyclic Cayley coordinates.
\end{abstract}


\section{Introduction}\label{extension:section:01}\label{sec:introduction}

Polyphase decompositions, filter banks, lifting schemes, and wavelet packets are standard parts of multirate signal processing and harmonic analysis. Their classical development is described in \cite{CrochiereRabiner1983,Vaidyanathan1993,VetterliKovacevic1995,StrangNguyen1996,Daubechies1992,Mallat2009}. These theories are often introduced through several parallel languages. Residue classes appear in polyphase analysis. Matrix symbols appear in filter banks. Trees appear in wavelet packets. Singularities are handled through a different approximation or detection method. The first aim of this article is to place these constructions in one exact coordinate system.

The second aim is to separate representation from recovery. If a signal has already been written as
\[
f=s+j
\]
with a smooth part \(s\) and a singular packet \(j\), then
\[
\widetilde f=Qs+j
\]
satisfies
\[
f-\widetilde f=s-Qs.
\]
This identity is exact, but it does not discover \(j\). A useful signal theory must therefore distinguish three errors. The first is the approximation error of the smooth component. The second is the error in the estimated singular packet. The third is the amplification caused by applying a reconstruction operator to that packet error. The middle sections of the paper isolate all three terms.

The singular packets considered here range from indicator functions of regions to truncated powers
\[
\frac{(t-\tau)_+^m}{m!}.
\]
The latter encode jumps in derivatives of order \(m\). They permit a piecewise smooth function to be reduced to a globally Sobolev remainder under explicit trace assumptions. This gives a precise version of the informal statement that a discontinuity should be represented directly rather than approximated by a global smooth basis.

The recovery theory begins with one dimensional separated jumps. For bounded deterministic noise, thresholding adjacent sample contrasts recovers the exact grid cells containing the jumps whenever the jump size exceeds the combined smooth variation and noise by an explicit margin. The same proof gives location and amplitude bounds. A sub Gaussian version follows from a union bound. Averaging repeated measurements reduces the required threshold in the expected way. Haar coefficients provide a second detection mechanism. Two grids shifted by half a cell guarantee that every jump is well placed in at least one grid, which removes the most elementary dyadic alignment failure.

The literature on singularity detection is much broader than the results proved here. Wavelet modulus maxima are developed in \cite{MallatHwang1992}. Best basis selection in packet libraries was introduced in \cite{CoifmanWickerhauser1992}. Nonlinear approximation is surveyed in \cite{DeVore1998}. Geometric image models, wedgelets, and curvelets appear in \cite{Donoho1999,CandesDonoho2004}. The present article does not replace those theories. Its narrower purpose is to give a common exact calculus in which coordinate reconstruction, adaptive basis selection, singular packet estimation, and smooth approximation remain logically separate.

The article is organized as follows. \Cref{sec:signal_seeds,sec:polyphase} develop lattice addresses and exact residue coordinates in any dimension. \Cref{sec:convolution,sec:filter_banks} express convolution and filter banks in those coordinates, including stability under imperfect reconstruction and filter perturbation. \Cref{sec:packet_trees} treats arbitrary pruned packet trees and stability of the best basis decision. \Cref{sec:singular_packets,sec:packet_error} introduce step and higher order singular packets, exact reconstruction identities, and error formulas. \Cref{sec:jump_recovery,sec:localization} prove deterministic and probabilistic recovery results. \Cref{sec:gibbs,sec:nonlinear} treat Gibbs oscillation and nonlinear interactions. The final sections record finite numerical illustrations and a secondary cyclic Cayley realization.

Several conclusions are deliberately excluded. The paper does not prove a universal image segmentation theorem. Exact recovery of separated jumps in one dimension does not imply stable recovery of arbitrary boundaries in higher dimensions. Retaining a correct singular packet does not remove noise or model error. An adaptive tree selected by an additive objective need not remain optimal when global hardware costs or cross packet interactions are introduced. The Cayley coordinate section is a secondary conformal realization and is not used in the recovery theorems.

\subsection{Main results}

The first result is an exact coordinate theorem. Let \(L\) be a finite index sublattice of \(\ZZ^d\), and let \(R\) contain one representative from every coset. The analysis map
\[
x\longmapsto (x_r)_{r\in R},
\qquad
x_r(\ell)=x(r+\ell),
\]
is an isometric isomorphism on every \(\ell^p\) space. This statement includes the usual one dimensional polyphase decomposition, but it also applies to arbitrary finite index lattice packets in several dimensions. If \(L'\leq L\), the coarse packets split exactly into their fine descendants. A convolution operator becomes a matrix of convolution operators on \(L\), with entries
\[
h_{rs}(m)=h(r-s+m).
\]
The matrix is therefore a coordinate representation of the original operator, not a reduced model.

The filter bank results distinguish exact reconstruction from quantitative near reconstruction. Exact recovery with delay \(d\) is the pointwise matrix identity
\[
R(z)E(z)=z^{-d}I.
\]
When this identity fails, the reconstruction error is controlled by the essential supremum of the delay compensated multiplier defect. If the filters themselves are perturbed, \cref{thm:filter_perturbation} separates the analysis error, the synthesis error, and their product. The frame bound in \cref{thm:frame_perturbation} shows that a lower singular value gap survives every sufficiently small perturbation. This makes the distinction between algebraic invertibility and stable numerical reconstruction explicit.

The packet tree results treat every admissible pruning, not only complete trees of a fixed depth. Unitary local splits give a unitary global transform. Local frame bounds propagate recursively in the nonunitary case. If the cost of a pruning is additive over retained nodes, then
\[
F(v)
=
\min\left\{
c(v),
\sum_{w\in\operatorname{child}(v)}F(w)
\right\}
\]
gives the exact optimal tree. The cost and the selected basis are stable when the node costs change by less than the gap separating the optimal pruning from its competitors.

The singular packet results begin with trace data. If a piecewise smooth function has a jump in its derivative of order \(m\) at \(\tau_q\), then the packet
\[
\jump{f^{(m)}}_{\tau_q}
\frac{(t-\tau_q)_+^m}{m!}
\]
removes exactly that trace jump. Subtracting all packets through order \(r-1\) leaves a global \(W^{r,p}\) remainder. The packet coefficients are uniquely determined by the one sided trace jumps. This gives a finite and explicit singular part for the class covered by \cref{thm:higher_order_packet_decomposition}.

The reconstruction formula with an estimated packet is
\[
f-\widehat f_m
=
(I-Q_m)s+(I-Q_m)(j-\widehat j).
\]
This identity identifies the two terms that must be controlled. The exact packet formulas in \cref{sec:packet_error} then evaluate or bound \(\|j-\widehat j\|\) from amplitude error, region overlap, boundary displacement, and the order of a derivative jump. A value jump has location modulus \(|\tau-\widehat\tau|^{1/p}\) in finite \(L^p\), while every positive order truncated power packet is Lipschitz in its location.

The principal recovery theorem uses adjacent sample contrasts. Let
\[
b_h=Mh+2\varepsilon
\]
be the largest contrast that can be produced by a Lipschitz background and bounded noise in one grid cell. If
\[
a_{\min}>2b_h,
\]
then every threshold between \(b_h\) and \(a_{\min}-b_h\) detects exactly the cells containing jumps. The jump locations are recovered within \(h/2\), and each amplitude error is at most \(b_h\). The probability theorem replaces the deterministic noise term by
\[
2\sigma\sqrt{\log(2n/\delta)}.
\]
Two half shifted Haar grids provide a second recovery mechanism. Every point is at least one quarter of a cell away from the boundary of one of the grids, so the Haar coefficient gap can be made independent of dyadic alignment.

\subsection{Notation}

The following symbols recur throughout the article.
\begin{center}
\begin{tabular}{p{0.22\textwidth}p{0.66\textwidth}}
\toprule
symbol & meaning\\
\midrule
\(L\leq\ZZ^d\) & a finite index sublattice\\
\(R\) & a set of representatives for \(\ZZ^d/L\)\\
\(P_r\) & the coordinate projector onto the coset \(r+L\)\\
\(\Acal_L,\Scal_L\) & lattice analysis and interleaving synthesis\\
\(E,R\) & filter bank analysis and synthesis symbols\\
\(Q_m\) & a bounded linear approximation operator acting on the smooth part\\
\(j,\widehat j\) & the true and estimated singular packets\\
\(H_\tau\) & the step packet \(\one_{[\tau,1]}\)\\
\(\phi_{m,\tau}\) & the truncated power packet \((t-\tau)_+^m/m!\)\\
\(h\) & the grid spacing in the recovery theorems\\
\bottomrule
\end{tabular}
\end{center}

\section{Signal seeds and lattice addresses}\label{extension:section:02}\label{sec:signal_seeds}

Let \(1\leq p\leq\infty\). For a countable set \(G\), write \(\ell^p(G)\) for the usual complex sequence space. For an integer \(d\geq1\), when \(G=\ZZ^d\), translation by \(a\in\ZZ^d\) is
\[
(T_ax)(n)=x(n-a).
\]
Let \(\delta_0\) be the impulse at the origin.

\begin{theorem}[Impulse generation]\label{thm:impulse_generation}
For \(1\leq p<\infty\),
\[
\overline{\Span\{T_a\delta_0:a\in\ZZ^d\}}^{\ell^p}
=
\ell^p(\ZZ^d).
\]
The same statement holds in \(c_0(\ZZ^d)\) with the supremum norm. In \(\ell^\infty(\ZZ^d)\), the closed span of the translation orbit is exactly \(c_0(\ZZ^d)\).
\end{theorem}

\begin{proof}
The translation orbit is the standard coordinate family \((\delta_a)_{a\in\ZZ^d}\). Its finite span is the space of finitely supported sequences. Finitely supported sequences are dense in \(\ell^p\) for finite \(p\) and in \(c_0\). Their closure in the supremum norm is \(c_0\), which is a proper closed subspace of \(\ell^\infty\).
\end{proof}

\begin{remark}\label{rem:infinity_boundary}
The distinction at \(p=\infty\) is structural. The impulse generates every finitely supported signal algebraically, but it is not a topological cyclic seed for all bounded signals.
\end{remark}

\begin{proposition}[Finite periodic signals]\label{prop:finite_periodic}
Let \(G=\ZZ^d/L\), where \(L\leq\ZZ^d\) has finite index. The translation orbit of the impulse at the identity is a basis of \(\CC^G\). In particular, one impulse generates every signal on the finite periodic lattice.
\end{proposition}

\begin{proof}
Translation sends the impulse at the identity to the impulse at each element of \(G\). These impulses form the standard basis of \(\CC^G\).
\end{proof}

\section{Polyphase coordinates on finite index sublattices}\label{extension:section:03}\label{sec:polyphase}

Let \(L\leq\ZZ^d\) be a sublattice of finite index \(q\). Choose a set \(R\subset\ZZ^d\) containing one representative from every coset in \(\ZZ^d/L\). Thus
\[
\ZZ^d=\bigsqcup_{r\in R}(r+L).
\]
For \(x:\ZZ^d\to\CC\), its channel at \(r\) is the sequence on \(L\)
\[
x_r(\ell)=x(r+\ell).
\]

\begin{definition}[Lattice packet]\label{def:lattice_packet}
For \(r\in R\), define the coordinate projector
\[
(P_rx)(n)=\one_{r+L}(n)x(n).
\]
The family \((P_r)_{r\in R}\) is called the packet system associated with \(L\) and \(R\).
\end{definition}

For finite \(p\), equip \(\bigoplus_{r\in R}\ell^p(L)\) with
\[
\|(u_r)\|^p=\sum_{r\in R}\|u_r\|_p^p.
\]
For \(p=\infty\), use the maximum norm.

\begin{theorem}[Exact lattice polyphase decomposition]\label{thm:lattice_polyphase}
For every \(1\leq p\leq\infty\), the analysis map
\[
\Acal_L:\ell^p(\ZZ^d)\longrightarrow\bigoplus_{r\in R}\ell^p(L),
\qquad
\Acal_Lx=(x_r)_{r\in R},
\]
is an isometric isomorphism. Its inverse is the interleaving map
\[
\Scal_L(u)(r+\ell)=u_r(\ell).
\]
The projectors satisfy
\[
P_rP_s=\delta_{rs}P_r,
\qquad
\sum_{r\in R}P_r=I.
\]
\end{theorem}

\begin{proof}
Every address has a unique representation \(r+\ell\). For finite \(p\),
\[
\|x\|_p^p
=
\sum_{r\in R}\sum_{\ell\in L}|x(r+\ell)|^p
=
\sum_{r\in R}\|x_r\|_p^p.
\]
The corresponding supremum identity is immediate when \(p=\infty\). Unique address representation proves that analysis and interleaving are inverse maps. The projector identities follow from the disjoint coset partition.
\end{proof}

\begin{corollary}[One dimensional residue coordinates]\label{cor:one_dimensional_polyphase}
For \(L=N\ZZ\) and \(R=\{0,\ldots,N-1\}\),
\[
x_r(k)=x(Nk+r)
\]
and
\[
\ell^p(\ZZ)\cong\bigoplus_{r=0}^{N-1}\ell^p(\ZZ)
\]
isometrically.
\end{corollary}

\begin{proof}
Apply \cref{thm:lattice_polyphase} to the standard residue representatives.
\end{proof}

\begin{proposition}[Translation covariance]\label{prop:translation_covariance}
Fix \(a\in\ZZ^d\). For every \(r\in R\), write uniquely
\[
r-a=r_a+\lambda_{a,r},
\qquad
r_a\in R,
\quad
\lambda_{a,r}\in L.
\]
Then
\[
(\Acal_LT_ax)_r(\ell)
=
x_{r_a}(\ell+\lambda_{a,r}).
\]
Thus a fine translation permutes the lattice packets and applies an internal translation on \(L\).
\end{proposition}

\begin{proof}
Evaluate at \(r+\ell\):
\[
(T_ax)(r+\ell)
=x(r+\ell-a)
=x(r_a+\ell+\lambda_{a,r}).
\]
This is the stated channel formula.
\end{proof}

\begin{figure}
\centering
\includegraphics[width=0.82\textwidth]{01_polyphase_coordinates.pdf}
\caption{A signal and its four exact residue channels. Interleaving the channels returns the original samples without approximation.}
\label{fig:polyphase_coordinates}
\end{figure}

\subsection{Refinement of sublattices}

Let \(L'\leq L\leq\ZZ^d\), with both indices finite. Choose representative systems \(R\) for \(\ZZ^d/L\) and \(R'\) for \(\ZZ^d/L'\). For \(r\in R\), define
\[
\operatorname{Child}(r)
=
\{s\in R':s+L'\subset r+L\}.
\]

\begin{theorem}[Exact lattice refinement]\label{thm:lattice_refinement}
For every coarse packet,
\[
P_r^{L}
=
\sum_{s\in\operatorname{Child}(r)}P_s^{L'}.
\]
The fine packets in the sum are pairwise disjoint. Refinement through several nested sublattices is transitive.
\end{theorem}

\begin{proof}
The coset \(r+L\) is the disjoint union of the \(L'\) cosets contained in it. Multiplying by the corresponding indicator functions gives the projector identity. A nested partition of a coset is independent of whether the intermediate level is inserted, which proves transitivity.
\end{proof}

\begin{example}[Rectangular image phases]\label{ex:image_phases}
For
\[
L=M\ZZ\times N\ZZ\subset\ZZ^2,
\]
one may take
\[
R=\{0,\ldots,M-1\}\times\{0,\ldots,N-1\}.
\]
The analysis map separates an image into \(MN\) exact phases. Refining \((M,N)\) to \((M',N')\), where \(M\mid M'\) and \(N\mid N'\), splits every coarse phase into
\[
\frac{M'}M\frac{N'}N
\]
fine phases.
\end{example}

\section{Convolution as a packet transfer matrix}\label{extension:section:04}\label{sec:convolution}

Let \(h\in\ell^1(\ZZ^d)\) and \(x\in\ell^p(\ZZ^d)\), with \(1\leq p\leq\infty\), and define
\[
(Hx)(n)
=(h*x)(n)
=\sum_{q\in\ZZ^d}h(q)x(n-q).
\]
Young's inequality gives
\[
\|Hx\|_p\leq\|h\|_1\|x\|_p
\]
for every \(1\leq p\leq\infty\).

\begin{definition}[Packet transfer kernel]\label{def:packet_transfer_kernel}
For output residue \(r\in R\) and input residue \(s\in R\), define a kernel on \(L\) by
\[
h_{rs}(m)=h(r-s+m),
\qquad
m\in L.
\]
\end{definition}

\begin{theorem}[Exact convolution transfer law]\label{thm:convolution_transfer}
Let \(y=h*x\). Then
\[
y_r(\ell)
=
\sum_{s\in R}\sum_{m\in L}h_{rs}(m)x_s(\ell-m).
\]
Equivalently,
\[
\Acal_LH\Acal_L^{-1}
=
(H_{rs})_{r,s\in R},
\]
where \(H_{rs}\) is convolution on \(L\) by \(h_{rs}\).
\end{theorem}

\begin{proof}
Evaluate the output at \(r+\ell\). Write every input address as \(s+\ell'\). The convolution offset is
\[
(r+\ell)-(s+\ell')
=(r-s)+(\ell-\ell').
\]
Setting \(m=\ell-\ell'\) gives the formula.
\end{proof}

\begin{proposition}[Transfer norm control]\label{prop:transfer_norm_control}
For every output residue \(r\),
\[
\sum_{s\in R}\|h_{rs}\|_{\ell^1(L)}
=
\|h\|_{\ell^1(\ZZ^d)}.
\]
The packet transfer matrix acts on the polyphase direct sum with norm at most \(\|h\|_1\) on every \(\ell^p\) space.
\end{proposition}

\begin{proof}
For fixed \(r\), every \(q\in\ZZ^d\) has a unique representation \(q=r-s+m\) with \(s\in R\) and \(m\in L\). This proves the equality of the absolute sums. The operator norm bound also follows by conjugating Young's inequality through the isometric analysis map.
\end{proof}

\begin{proposition}[Composition of transfer matrices]\label{prop:transfer_composition}
If \(G\) and \(H\) are convolution operators with packet matrices \((G_{rs})\) and \((H_{rs})\), then the packet matrix of \(GH\) is
\[
(GH)_{rt}
=
\sum_{s\in R}G_{rs}H_{st}.
\]
In one dimensional transform notation, the corresponding matrix symbols multiply pointwise.
\end{proposition}

\begin{proof}
Conjugation by \(\Acal_L\) preserves composition:
\[
\Acal_LGH\Acal_L^{-1}
=
(\Acal_LG\Acal_L^{-1})(\Acal_LH\Acal_L^{-1}).
\]
Reading the \((r,t)\) block gives the formula.
\end{proof}

\subsection{The one dimensional symbol}

For \(L=N\ZZ\) with \(N\geq1\), first take \(|z|=1\). The assumption \(h\in\ell^1(\ZZ)\) makes the following symbol series absolutely convergent. Define
\[
H_{ks}(z)=\sum_{m\in\ZZ}h(Nm+k-s)z^{-m}.
\]
For \(x\in\ell^1(\ZZ)\), the channel transforms converge absolutely on the unit circle and satisfy
\[
\bm Y(z)=\bm H_N(z)\bm X(z).
\]
For \(x\in\ell^2(\ZZ)\), the same identity holds almost everywhere in the \(L^2\) Fourier sense. For finite impulse-response kernels it also holds at nonzero \(z\) wherever the input channel transforms are weighted absolutely summable. This is the classical polyphase matrix. It is the exact block matrix of the original convolution operator.

\begin{corollary}[Finite impulse response closure]\label{cor:fir_closure}
If \(h\) has finite support, every packet kernel \(h_{ks}\) has finite support. Matrix multiplication of polyphase symbols agrees with composition of the original finite impulse response operators.
\end{corollary}

\begin{proof}
Finite support is preserved by taking residue subsequences. The composition statement is \cref{prop:transfer_composition}.
\end{proof}

\begin{example}[A four phase finite impulse response operator]\label{ex:four_phase_fir}
Let
\[
h(0)=1,
\quad h(1)=2,
\quad h(2)=3,
\quad h(3)=4,
\quad h(4)=5,
\]
with all other taps zero. For four residue channels, the transfer symbol is
\[
\bm H_4(z)
=
\begin{pmatrix}
1+5z^{-1}&4z^{-1}&3z^{-1}&2z^{-1}\\
2&1+5z^{-1}&4z^{-1}&3z^{-1}\\
3&2&1+5z^{-1}&4z^{-1}\\
4&3&2&1+5z^{-1}
\end{pmatrix}.
\]
The delay term records the tap that crosses a coarse cell boundary. No approximation is introduced by the matrix representation.
\end{example}

\begin{figure}
\centering
\includegraphics[width=0.62\textwidth]{02_transfer_blocks.pdf}
\caption{Absolute tap weight in the transfer blocks of the four phase operator in \cref{ex:four_phase_fir}.}
\label{fig:transfer_blocks}
\end{figure}

\section{Filter banks, frame bounds, and lifting}\label{extension:section:05}\label{sec:filter_banks}

Throughout this section the banks act boundedly on the relevant polyphase \(\ell^2\) spaces. A finite channel analysis bank acts in Fourier coordinates by a matrix \(E(z)\). A synthesis bank acts by a matrix \(R(z)\). The matrices may be square or rectangular. Their entries are Laurent polynomials for finite filters and bounded measurable functions for general stable filters.

\begin{theorem}[Perfect reconstruction criterion]\label{thm:perfect_reconstruction}
Let \(E(z)\) be an \(M\) by \(N\) analysis matrix and let \(R(z)\) be an \(N\) by \(M\) synthesis matrix. The bank reconstructs every input with a delay of \(d\) coarse channel samples, equivalent to \(Nd\) original samples in the one dimensional \(N\)-phase convention, exactly if and only if
\[
R(z)E(z)=z^{-d}I_N
\]
for almost every point on the unit circle. A nonzero scalar factor gives exact reconstruction up to gain and delay.
\end{theorem}

\begin{proof}
The channel Fourier transform turns analysis followed by synthesis into multiplication by \(R(z)E(z)\). If the matrix identity holds, inverse transformation gives the delayed input. Conversely, if every input is reconstructed with that delay, the difference matrix annihilates every vector valued Fourier datum and must vanish almost everywhere.
\end{proof}

\begin{theorem}[Quantitative imperfect reconstruction]\label{thm:imperfect_reconstruction}
Define the delay compensated defect
\[
D(\e^{\ii\omega})
=
\e^{\ii d\omega}R(\e^{\ii\omega})E(\e^{\ii\omega})-I_N.
\]
If
\[
\varepsilon_{\mathrm{rec}}
=
\esssup_{\omega}\|D(\e^{\ii\omega})\|_{\mathrm{op}}<\infty,
\]
then the delay compensated reconstruction \(\widetilde x\) satisfies
\[
\|\widetilde x-x\|_2
\leq
\varepsilon_{\mathrm{rec}}\|x\|_2.
\]
The constant is sharp as an operator norm.
\end{theorem}

\begin{proof}
In polyphase Fourier coordinates, the error is \(D(\e^{\ii\omega})\bm X(\e^{\ii\omega})\). Pointwise multiplication and Plancherel's theorem give the upper bound. The operator norm of a measurable matrix multiplier is its essential supremum, which proves sharpness.
\end{proof}

\begin{example}[A reconstruction certificate]\label{ex:reconstruction_certificate}
Suppose a two channel bank has delay compensated defect
\[
D(\e^{\ii\omega})
=
\begin{pmatrix}
0.01\cos\omega&0\\
0&0.02\sin\omega
\end{pmatrix}.
\]
Then
\[
\esssup_\omega\|D(\e^{\ii\omega})\|_{\mathrm{op}}
=0.02.
\]
Every input therefore has relative reconstruction error at most two percent in \(\ell^2\). The estimate is independent of the input length.
\end{example}

\begin{theorem}[Stability under filter perturbation]\label{thm:filter_perturbation}
Assume \(RE=z^{-d}I_N\). Let
\[
\widetilde E=E+\Delta E,
\qquad
\widetilde R=R+\Delta R.
\]
Set
\[
\begin{aligned}
e_0&=\esssup\|E\|_{\mathrm{op}},
& r_0&=\esssup\|R\|_{\mathrm{op}},\\
\delta_E&=\esssup\|\Delta E\|_{\mathrm{op}},
& \delta_R&=\esssup\|\Delta R\|_{\mathrm{op}}.
\end{aligned}
\]
Then the delay compensated reconstruction defect obeys
\[
\varepsilon_{\mathrm{rec}}
\leq
r_0\delta_E+e_0\delta_R+\delta_R\delta_E.
\]
\end{theorem}

\begin{proof}
Expand
\[
\widetilde R\widetilde E-RE
=
R\Delta E+\Delta R E+\Delta R\Delta E.
\]
Take pointwise operator norms, then the essential supremum. Multiplication by the delay factor has unit modulus.
\end{proof}

\begin{theorem}[Frame bounds from the polyphase symbol]\label{thm:frame_bounds}
Assume that
\[
A I_N
\leq
E(\e^{\ii\omega})^*E(\e^{\ii\omega})
\leq
B I_N
\]
for almost every \(\omega\), where \(0<A\leq B<\infty\). Then the analysis bank satisfies
\[
A\|x\|_2^2
\leq
\|\Acal_E x\|_2^2
\leq
B\|x\|_2^2.
\]
The pointwise matrix
\[
R_{\rm can}(\e^{\ii\omega})
=
(E^*E)^{-1}E^*
\]
is a bounded left inverse and gives the canonical dual synthesis bank.
\end{theorem}

\begin{proof}
Apply the matrix inequality to the polyphase Fourier vector at every frequency and integrate. Plancherel's theorem gives the frame inequalities. The lower bound makes \(E^*E\) invertible almost everywhere, and direct multiplication gives \(R_{\rm can}E=I_N\). The same spectral bounds show that the canonical dual is bounded.
\end{proof}

\begin{corollary}[Coefficient noise under canonical synthesis]\label{cor:canonical_noise}
Under the hypotheses of \cref{thm:frame_bounds}, suppose the observed coefficient vector is
\[
c=\Acal_Ex+\nu.
\]
Canonical synthesis gives
\[
\widehat x=R_{\mathrm{can}}c
\]
and
\[
\|\widehat x-x\|_2
\leq
\frac1{\sqrt A}\|\nu\|_2.
\]
\end{corollary}

\begin{proof}
The smallest singular value of \(E\) is at least \(\sqrt A\), so the pointwise norm of its pseudoinverse is at most \(1/\sqrt A\). Apply Plancherel's theorem to \(R_{\mathrm{can}}\nu\).
\end{proof}

\begin{theorem}[Perturbation of frame bounds]\label{thm:frame_perturbation}
Under the hypotheses of \cref{thm:frame_bounds}, let \(\widetilde E=E+\Delta E\) and suppose
\[
\delta=\esssup_\omega\|\Delta E(\e^{\ii\omega})\|_{\mathrm{op}}<\sqrt A.
\]
Then \(\widetilde E\) is a frame analysis bank with bounds
\[
(\sqrt A-\delta)^2
\quad\text{and}\quad
(\sqrt B+\delta)^2.
\]
\end{theorem}

\begin{proof}
For every vector \(v\),
\[
\|\widetilde Ev\|
\geq
\|Ev\|-\|\Delta Ev\|
\geq
(\sqrt A-\delta)\|v\|.
\]
The upper bound follows from
\[
\|\widetilde Ev\|
\leq
(\sqrt B+\delta)\|v\|.
\]
Square the inequalities and integrate in frequency.
\end{proof}

\begin{remark}\label{rem:canonical_dual_support}
Even when \(E\) is a Laurent polynomial matrix, the canonical dual need not have finite impulse response. Finite support of both analysis and synthesis is an additional algebraic condition.
\end{remark}

\begin{corollary}[Paraunitary bank]\label{cor:paraunitary}
If \(E(z)^*E(z)=I_N\) on the unit circle, then analysis is an isometry and synthesis may be chosen as \(R(z)=E(z)^*\).
\end{corollary}

\begin{proof}
Use \cref{thm:frame_bounds} with \(A=B=1\), then multiply the matrices.
\end{proof}

\begin{figure}
\centering
\includegraphics[width=0.72\textwidth]{03_frame_stability.pdf}
\caption{Pointwise singular value bounds govern both the original frame and its perturbed analysis bank.}
\label{fig:frame_stability}
\end{figure}

\subsection{Haar analysis}

The Haar matrix is
\[
E_H
=
\frac1{\sqrt2}
\begin{pmatrix}
1&1\\
1&-1
\end{pmatrix}.
\]
For even and odd samples \(e(n)=x(2n)\) and \(o(n)=x(2n+1)\), the low and high coordinates are
\[
a(n)=\frac{e(n)+o(n)}{\sqrt2},
\qquad
d(n)=\frac{e(n)-o(n)}{\sqrt2}.
\]

\begin{corollary}[Haar reconstruction and energy]\label{cor:haar_reconstruction}
The Haar bank reconstructs exactly and
\[
\|x\|_2^2=\|a\|_2^2+\|d\|_2^2.
\]
\end{corollary}

\begin{proof}
The matrix \(E_H\) is unitary.
\end{proof}

\subsection{Lifting steps}

Let \(P(z)\) and \(Q(z)\) be Laurent polynomials. Define
\[
L_P(z)
=
\begin{pmatrix}
1&0\\
P(z)&1
\end{pmatrix},
\qquad
U_Q(z)
=
\begin{pmatrix}
1&Q(z)\\
0&1
\end{pmatrix}.
\]

\begin{theorem}[Exact lifting update]\label{thm:lifting_update}
The lifting matrices are invertible with
\[
L_P^{-1}=L_{-P},
\qquad
U_Q^{-1}=U_{-Q}.
\]
If \(R_0E_0=I\), then every finite product \(V\) of lifting matrices gives another exact bank by
\[
E=VE_0,
\qquad
R=R_0V^{-1}.
\]
Every lifting step has determinant one.
\end{theorem}

\begin{proof}
Direct multiplication proves the inverse formulas and the determinant statement. The reconstructed product is
\[
RE
=R_0V^{-1}VE_0
=R_0E_0
=I.
\]
\end{proof}

\begin{example}[A two neighbor predict and update transform]\label{ex:predict_update}
Split a sequence into even and odd samples. Define
\[
d(n)
=o(n)-\frac{e(n)+e(n+1)}2
\]
and
\[
s(n)
=e(n)+\frac{d(n-1)+d(n)}4.
\]
The inverse is obtained in reverse order:
\[
e(n)
=s(n)-\frac{d(n-1)+d(n)}4,
\]
then
\[
o(n)
=d(n)+\frac{e(n)+e(n+1)}2.
\]
Thus reconstruction is exact over every coefficient field in which two is invertible. This is a standard compactly supported biorthogonal lifting pattern \cite{CohenDaubechiesFeauveau1992,Sweldens1996,DaubechiesSweldens1998}.
\end{example}

\section{Wavelet packet trees and admissible bases}\label{extension:section:06}\label{sec:packet_trees}

A packet transform is built by deciding at each node whether to retain the current coordinate space or to split it into children. The proof of reconstruction is local and does not depend on a particular wavelet family.

\begin{definition}[Unitary packet tree]\label{def:unitary_packet_tree}
A finite unitary packet tree consists of Hilbert spaces \(\Hcal_v\) at its vertices and, at every internal vertex \(v\), a unitary map
\[
U_v:\Hcal_v\longrightarrow\bigoplus_{w\in\operatorname{child}(v)}\Hcal_w.
\]
An admissible pruning is a set of vertices that meets every root to terminal path exactly once.
\end{definition}

\begin{theorem}[Reconstruction on a pruned packet tree]\label{thm:pruned_packet_tree}
Every admissible pruning \(\Lcal\) defines a unitary analysis map
\[
\Acal_{\Lcal}:\Hcal_{\rm root}
\longrightarrow
\bigoplus_{v\in\Lcal}\Hcal_v.
\]
Its inverse is obtained by applying the adjoints of the local splits from the leaves toward the root. In particular,
\[
\|x\|^2
=
\sum_{v\in\Lcal}\|x_v\|^2.
\]
\end{theorem}

\begin{proof}
Proceed by induction on the depth of the retained subtree. If the root is retained, the statement is immediate. If the root is split, apply the unitary map \(U_{\rm root}\), then apply the induction hypothesis independently in each child subtree. A direct sum of unitary maps is unitary. Reversing the maps with adjoints gives synthesis.
\end{proof}

\begin{corollary}[Coefficient perturbations in a unitary tree]\label{cor:unitary_tree_noise}
Let \(c=\Acal_{\Lcal}x\) be the leaf coefficients of an admissible unitary pruning, and let \(\widehat c=c+e\). If \(\widehat x\) is synthesized from \(\widehat c\), then
\[
\|\widehat x-x\|=\|e\|.
\]
Thus coefficient quantization and coefficient noise are neither amplified nor attenuated by exact unitary synthesis.
\end{corollary}

\begin{proof}
The synthesis map is unitary by \cref{thm:pruned_packet_tree}. Hence
\[
\widehat x-x
=\Acal_{\Lcal}^{-1}e
\]
has the same norm as \(e\).
\end{proof}

\begin{proposition}[Bounds for nonunitary local splits]\label{prop:nonunitary_tree_bounds}
Suppose every internal node has an analysis map \(U_v\) and constants \(0\leq A_v\leq B_v<\infty\) satisfying
\[
A_v\|x\|^2
\leq
\sum_{w\in\operatorname{child}(v)}\|(U_vx)_w\|^2
\leq
B_v\|x\|^2.
\]
Define recursively
\[
a(v)=
\begin{cases}
1,&v\text{ retained},\\
A_v\min_w a(w),&v\text{ split},
\end{cases}
\]
and
\[
b(v)=
\begin{cases}
1,&v\text{ retained},\\
B_v\max_w b(w),&v\text{ split}.
\end{cases}
\]
Then the leaf analysis of the pruned tree has bounds \(a({\rm root})\) and \(b({\rm root})\).
\end{proposition}

\begin{proof}
Assume the child subtree bounds have been proved. Then
\[
\sum_{w}\sum_{\ell\in\Lcal_w}\|x_\ell\|^2
\geq
\min_w a(w)
\sum_w\|(U_vx)_w\|^2
\geq
A_v\min_w a(w)\|x\|^2.
\]
The upper bound follows in the same way with the maximum and \(B_v\). Induction completes the proof.
\end{proof}

\begin{figure}
\centering
\includegraphics[width=0.72\textwidth]{04_packet_tree.pdf}
\caption{A complete packet tree and one admissible pruning. Retained nodes may occur at different depths.}
\label{fig:packet_tree}
\end{figure}

\subsection{Best basis recursion}

\begin{interpretation}[title={The minimum is taken over feasible prunings}]
The value at a node is the least cost among all admissible prunings of the subtree rooted there. It is one optimal value. Individual competing prunings may have larger costs. The recursion compares retaining the node with the sum of the already optimal child values, with retention chosen on a tie.
\end{interpretation}

Assign a real cost \(c(v)\) to retaining each node. For an admissible pruning \(\Lcal\), let
\[
C(\Lcal)=\sum_{v\in\Lcal}c(v).
\]
The cost may measure entropy, coefficient count, an approximation penalty, or any other additive local objective.

\begin{theorem}[Exact additive best basis recursion]\label{thm:best_basis}
For each vertex define
\[
F(v)
=
\min\left\{
 c(v),
 \sum_{w\in\operatorname{child}(v)}F(w)
\right\},
\]
with \(F(v)=c(v)\) at terminal vertices. Then \(F(v)\) is the minimum cost of every admissible pruning of the subtree rooted at \(v\). A minimizing pruning is recovered by retaining \(v\) when the first term is no larger than the second and splitting it otherwise.
\end{theorem}

\begin{proof}
Every admissible pruning of the subtree has exactly two forms. It either retains \(v\), with cost \(c(v)\), or splits \(v\) and independently chooses an admissible pruning in every child subtree. The minimum in the second case is the sum of the child minima. This gives the recursion and the reconstruction rule. In an equality case both choices minimize, and the stated rule retains the node.
\end{proof}

\begin{theorem}[Stability of the optimal cost]\label{thm:best_basis_cost_stability}
Let \(c\) and \(\widetilde c\) be two node cost functions. For an admissible pruning \(\Lcal\), set
\[
\varepsilon(\Lcal)
=
\sum_{v\in\Lcal}|c(v)-\widetilde c(v)|
\]
and let
\[
\varepsilon_*
=
\max_{\Lcal}\varepsilon(\Lcal).
\]
If \(F_*\) and \(\widetilde F_*\) are the two optimal root costs, then
\[
|F_*-\widetilde F_*|
\leq
\varepsilon_*.
\]
\end{theorem}

\begin{proof}
For every admissible pruning,
\[
|C(\Lcal)-\widetilde C(\Lcal)|
\leq
\varepsilon_*.
\]
Evaluate \(\widetilde C\) on a minimizer for \(C\) to obtain \(\widetilde F_*\leq F_*+\varepsilon_*\). Reverse the roles to obtain \(F_*\leq\widetilde F_*+\varepsilon_*\).
\end{proof}

\begin{corollary}[Stability of the selected basis]\label{cor:best_basis_selection_stability}
Assume the optimal pruning \(\Lcal_*\) for \(c\) is unique and at least one competing pruning exists. Define \(g=\min_{\Lcal\ne\Lcal_*}(C(\Lcal)-C(\Lcal_*))>0\). If
\[
2\varepsilon_*<g,
\]
then the optimal pruning for \(\widetilde c\) is the same. If only one admissible pruning exists, it remains selected under every cost change without a gap condition.
\end{corollary}

\begin{proof}
The cost of the original minimizer can increase by at most \(\varepsilon_*\), while the cost of every competitor can decrease by at most \(\varepsilon_*\). The original gap therefore remains positive.
\end{proof}

\begin{figure}
\centering
\includegraphics[width=0.72\textwidth]{05_best_basis.pdf}
\caption{An additive best basis recursion compares the cost of retaining a packet with the total cost of splitting it.}
\label{fig:best_basis}
\end{figure}

\section{Singular packets and exact reconstruction}\label{extension:section:07}\label{sec:singular_packets}

\begin{readerroadmap}[title={Trace data identify the singular part}]
A value jump needs a step packet. A jump in a derivative needs a truncated power of the corresponding order. Subtracting these finite packets matches the one sided traces and leaves the regular remainder covered by the theorem's Sobolev assumptions. The packet construction and the approximation operator play separate roles.
\end{readerroadmap}

A packet decomposition becomes useful for approximation when it separates a component that should be retained from a component that should be compressed. The simplest singular packets are steps. Higher order packets record jumps in derivatives.

For \(m\geq0\) and \(0<\tau<1\), define
\[
\phi_{m,\tau}(t)
=
\frac{(t-\tau)_+^m}{m!}.
\]
Thus \(\phi_{0,\tau}=\one_{[\tau,1]}\). On each side of \(\tau\), the packet is a polynomial. Its derivative of order \(m\) has a unit jump at \(\tau\).

\begin{definition}[Sobolev step packet class]\label{def:step_packet_class}
Fix \(r\geq1\), \(1\leq p\leq\infty\), and an integer \(Q\geq0\). The class \(\Fcal_{r,p}(Q)\) consists of functions
\[
f=s+\sum_{q=1}^{Q}a_q\phi_{0,\tau_q}
\]
with \(s\in W^{r,p}(0,1)\) and
\[
0<\tau_1<\cdots<\tau_Q<1.
\]
No lower bound on the jump amplitudes is required for representation. Recovery results later impose such a bound.
\end{definition}

\begin{theorem}[Higher order singular packet decomposition]\label{thm:higher_order_packet_decomposition}
Let
\[
0=\tau_0<\tau_1<\cdots<\tau_Q<\tau_{Q+1}=1.
\]
Suppose the restriction of \(f\) to every interval \((\tau_q,\tau_{q+1})\) belongs to \(W^{r,p}\), and choose the one dimensional Sobolev representatives whose derivatives through order \(r-1\) have one sided traces at the endpoints. Define
\[
a_{q,m}
=
\jump{f^{(m)}}_{\tau_q}
=
f^{(m)}(\tau_q+)-f^{(m)}(\tau_q-)
\]
and
\[
j_f(t)
=
\sum_{q=1}^{Q}\sum_{m=0}^{r-1}a_{q,m}\phi_{m,\tau_q}(t).
\]
Then
\[
s_f=f-j_f
\]
has matching one sided derivatives through order \(r-1\) at every breakpoint and belongs to \(W^{r,p}(0,1)\). Its weak derivative of order \(r\) equals the piecewise classical derivative \(f^{(r)}\) almost everywhere.
\end{theorem}

\begin{proof}
For \(0\leq \ell\leq m\), the classical derivative of \(\phi_{m,\tau}\) on the two open sides is
\[
\phi_{m,\tau}^{(\ell)}(t)
=
\frac{(t-\tau)_+^{m-\ell}}{(m-\ell)!}.
\]
This derivative has zero jump when \(\ell<m\) and a unit jump when \(\ell=m\). Derivatives of order larger than \(m\) vanish on both open sides. It follows that subtracting the term \(a_{q,m}\phi_{m,\tau_q}\) removes exactly the jump in the derivative of order \(m\) without changing the other classical trace jumps. Hence \(s_f\) has matching traces through order \(r-1\).

On each partition interval, \(j_f\) is a polynomial of degree at most \(r-1\), so its weak derivative of order \(r\) is zero there. A piecewise \(W^{r,p}\) function whose traces through order \(r-1\) match across every breakpoint belongs to \(W^{r,p}(0,1)\). Repeated integration by parts on the partition intervals shows that all boundary terms cancel. Thus the weak derivative of order \(r\) is the piecewise weak derivative of \(f\).
\end{proof}

\begin{corollary}[Uniqueness of the trace packet]\label{cor:trace_packet_uniqueness}
Fix the breakpoints \(\tau_q\). If
\[
f=s+\sum_{q=1}^{Q}\sum_{m=0}^{r-1}b_{q,m}\phi_{m,\tau_q}
\]
with \(s\in W^{r,p}(0,1)\), then
\[
b_{q,m}=\jump{f^{(m)}}_{\tau_q}.
\]
The finite truncated power packet is therefore the unique packet supported at the declared breakpoints that removes all trace jumps through order \(r-1\).
\end{corollary}

\begin{proof}
Every representative in \(W^{r,p}(0,1)\) has derivatives through order \(r-1\) that admit absolutely continuous representatives on compact subintervals, with matching one sided traces at each interior point. Taking the jump of the derivative of order \(m\) therefore removes the contribution of \(s\). The packet \(\phi_{m,\tau_q}\) has unit jump in its derivative of order \(m\), while all other declared packets have zero jump at that derivative order and breakpoint. Hence the remaining jump is exactly \(b_{q,m}\).
\end{proof}

\begin{remark}\label{rem:higher_order_scope}
The theorem isolates only finitely many trace jumps. More general free discontinuity models may contain jump sets, Cantor parts, or singular derivatives that are not represented by a finite truncated power sum. The natural ambient theory is related to bounded variation and free discontinuity spaces \cite{AmbrosioFuscoPallara2000}.
\end{remark}

\begin{figure}
\centering
\includegraphics[width=0.72\textwidth]{06_higher_order_packets.pdf}
\caption{Truncated power packets encode jumps in a function and in its successive derivatives.}
\label{fig:higher_order_packets}
\end{figure}

Let \(X\) be a normed linear space and let \(Q:X\to X\) be bounded and linear. The next identity is elementary, but it is the central separation principle.

\begin{theorem}[Exact retained packet identity]\label{thm:retained_packet_identity}
Let \(f=s+j\) in \(X\). Define
\[
\widetilde f=Qs+j.
\]
Then
\[
\boxed{
f-\widetilde f=(I-Q)s.
}
\]
Consequently,
\[
\|f-\widetilde f\|=\|s-Qs\|.
\]
\end{theorem}

\begin{proof}
Subtract:
\[
f-\widetilde f
=s+j-Qs-j
=(I-Q)s.
\]
\end{proof}

\begin{remark}\label{rem:retained_packet_boundary}
The theorem assumes that the packet \(j\) is known. It does not supply a detector. It also does not say that every representation of \(f\) as \(s+j\) is useful. The point is that once a meaningful packet has been identified, its approximation error can be removed from the smooth approximation problem exactly.
\end{remark}

\begin{corollary}[Sobolev rate after packet separation]\label{cor:sobolev_rate}
Suppose \(Q_m\) satisfies
\[
\|s-Q_ms\|_{L^p}
\leq
C2^{-mr}\|s\|_{W^{r,p}}
\]
for every \(s\in W^{r,p}(0,1)\). If \(f=s+j\) and \(j\) is retained, then
\[
\|f-(Q_ms+j)\|_{L^p}
\leq
C2^{-mr}\|s\|_{W^{r,p}}.
\]
The same conclusion applies to the higher order decomposition in \cref{thm:higher_order_packet_decomposition}.
\end{corollary}

\begin{proof}
Use \cref{thm:retained_packet_identity} and the assumed approximation estimate.
\end{proof}

\begin{theorem}[Reconstruction with an estimated packet]\label{thm:estimated_packet_reconstruction}
Let \(Q_m:X\to X\) be linear, let \(f=s+j\), let \(\widehat j\) be an estimated packet, and define
\[
\widehat f_m
=
Q_m(f-\widehat j)+\widehat j.
\]
Then
\[
\boxed{
f-\widehat f_m
=
(I-Q_m)s+(I-Q_m)(j-\widehat j).
}
\]
If \(Q_m\) is bounded, then
\[
\|f-\widehat f_m\|
\leq
\|s-Q_ms\|
+
(1+\|Q_m\|)\|j-\widehat j\|.
\]
\end{theorem}

\begin{proof}
Since \(f-\widehat j=s+(j-\widehat j)\),
\[
\widehat f_m
=
Q_ms+Q_m(j-\widehat j)+\widehat j.
\]
Subtract this expression from \(s+j\). The norm bound follows from
\[
\|(I-Q_m)u\|
\leq
(1+\|Q_m\|)\|u\|.
\]
\end{proof}

\begin{figure}
\centering
\includegraphics[width=0.84\textwidth]{07_known_packet_reconstruction.pdf}
\caption{A global Fourier approximation must represent the jumps. A retained packet approximation applies the Fourier truncation only to the smooth component.}
\label{fig:known_packet_reconstruction}
\end{figure}

\section{Exact error formulas for estimated packets}\label{extension:section:08}\label{sec:packet_error}

The error of an estimated packet is often described by a segmentation score. For approximation, an exact norm formula is more useful because it separates value errors from geometric disagreement.

\begin{theorem}[Exact scalar region formula]\label{thm:scalar_region_formula}
Let \(E,F\) be measurable subsets of a measure space, let \(a,b\in\CC\), and let \(1\leq p<\infty\). Then
\[
\begin{aligned}
\|a\one_E-b\one_F\|_p^p
={}&
|a-b|^p\mu(E\cap F)
+|a|^p\mu(E\setminus F)\\
&+|b|^p\mu(F\setminus E).
\end{aligned}
\]
For equal amplitudes,
\[
\|a\one_E-a\one_F\|_p
=
|a|\mu(E\mathbin\triangle F)^{1/p}.
\]
\end{theorem}

\begin{proof}
The sets \(E\cap F\), \(E\setminus F\), and \(F\setminus E\) are disjoint. The difference takes the values \(a-b\), \(a\), and \(-b\) on these three sets and is zero elsewhere. Integrating its \(p\)th power gives the formula.
\end{proof}

\begin{corollary}[Location error for one step]\label{cor:step_location_error}
Let \(0\leq\tau,\widehat\tau\leq1\) and \(H_\tau=\one_{[\tau,1]}\). For \(1\leq p<\infty\),
\[
\|aH_\tau-aH_{\widehat\tau}\|_p
=
|a||\tau-\widehat\tau|^{1/p}.
\]
In the supremum norm, the error equals \(|a|\) whenever \(\tau\neq\widehat\tau\).
\end{corollary}

\begin{proof}
The symmetric difference of the two intervals has length \(|\tau-\widehat\tau|\). The supremum statement follows because the two steps differ by \(a\) on a nonempty interval.
\end{proof}

\begin{theorem}[Location stability for higher order packets]\label{thm:higher_order_location}
Let \(m\geq1\), \(1\leq p<\infty\), and \(0\leq\tau,\widehat\tau\leq1\). Then
\[
\|\phi_{m,\tau}-\phi_{m,\widehat\tau}\|_{L^p(0,1)}
\leq
\frac{|\tau-\widehat\tau|}
{(m-1)!((m-1)p+1)^{1/p}}.
\]
For \(p=\infty\),
\[
\|\phi_{m,\tau}-\phi_{m,\widehat\tau}\|_\infty
\leq
\frac{|\tau-\widehat\tau|}{(m-1)!}.
\]
Thus location error is Lipschitz for every derivative jump packet of positive order, while a value jump has only the finite \(L^p\) modulus in \cref{cor:step_location_error}.
\end{theorem}

\begin{proof}
Assume \(\tau<\widehat\tau\). Differentiation with respect to the location parameter gives
\[
\frac{\partial}{\partial u}\phi_{m,u}(t)
=-\phi_{m-1,u}(t)
\]
for almost every \(t\). Hence
\[
\phi_{m,\tau}-\phi_{m,\widehat\tau}
=
\int_\tau^{\widehat\tau}\phi_{m-1,u}\dd u.
\]
Minkowski's inequality gives
\[
\|\phi_{m,\tau}-\phi_{m,\widehat\tau}\|_p
\leq
|\tau-\widehat\tau|
\sup_u\|\phi_{m-1,u}\|_p.
\]
The supremum is bounded by the norm of \(t^{m-1}/(m-1)!\) on \([0,1]\), which is the stated constant. The supremum norm argument is identical.
\end{proof}

\subsection{Several regions}

Let \((E_i)_{i=0}^{I}\) and \((F_j)_{j=0}^{J}\) be measurable partitions of the same domain, up to null sets. Let \(a_i\) and \(b_j\) be the corresponding values. The index zero may be used for a background region.

\begin{theorem}[Overlap matrix formula]\label{thm:overlap_matrix}
Define
\[
M_{ij}=\mu(E_i\cap F_j).
\]
For \(1\leq p<\infty\),
\[
\left\|
\sum_{i=0}^{I}a_i\one_{E_i}
-
\sum_{j=0}^{J}b_j\one_{F_j}
\right\|_p^p
=
\sum_{i=0}^{I}\sum_{j=0}^{J}|a_i-b_j|^pM_{ij}.
\]
The row sums of \(M\) are \(\mu(E_i)\), and the column sums are \(\mu(F_j)\).
\end{theorem}

\begin{proof}
The intersections \(E_i\cap F_j\) form a measurable partition. On such an intersection, the two piecewise constant functions have values \(a_i\) and \(b_j\). Integrate the constant error over every intersection. The marginal identities follow by summing the partition pieces in a fixed row or column.
\end{proof}

\begin{corollary}[Background and packet error]\label{cor:background_packet_error}
Suppose the true and estimated packets use the same amplitude list \((a_i)_{i=0}^{I}\), but possibly different regions. Then
\[
\|j-\widehat j\|_p^p
=
\sum_{i,k}|a_i-a_k|^p\mu(E_i\cap F_k).
\]
Only off diagonal overlaps contribute.
\end{corollary}

\begin{proof}
Apply \cref{thm:overlap_matrix} with \(b_k=a_k\). The diagonal terms vanish.
\end{proof}

\begin{figure}
\centering
\includegraphics[width=0.60\textwidth]{08_overlap_matrix.pdf}
\caption{An overlap matrix records both correct regions and every type of misclassification. The exact packet norm weights each entry by the corresponding value difference.}
\label{fig:overlap_matrix}
\end{figure}

\subsection{Boundary tube control}

For a closed set \(\Gamma\), write
\[
N_\delta(\Gamma)
=
\{x:\dist(x,\Gamma)<\delta\}.
\]

\begin{definition}[Linear tube bound]\label{def:linear_tube_bound}
A boundary \(\Gamma\) has a linear tube bound if there are constants \(C_\Gamma\) and \(\delta_0>0\) such that
\[
\mu(N_\delta(\Gamma))
\leq
C_\Gamma\delta
\]
for \(0<\delta<\delta_0\).
\end{definition}

Bounded Lipschitz domains and many finite perimeter sets have estimates of this form, with constants controlled by geometric size. The next theorem only requires the displayed bound.

\begin{theorem}[Geometric packet bound]\label{thm:geometric_packet_bound}
Let \(E\) be a region with boundary \(\Gamma\), and suppose
\[
E\mathbin\triangle F
\subseteq
N_\delta(\Gamma).
\]
If \(\Gamma\) has a linear tube bound and \(0<\delta<\delta_0\) lies in its stated range, then for \(1\leq p<\infty\),
\[
\|a\one_E-a\one_F\|_p
\leq
|a|(C_\Gamma\delta)^{1/p}.
\]
If amplitudes \(a\) and \(b\) differ, then
\[
\|a\one_E-b\one_F\|_p
\leq
|a-b|\mu(E\cap F)^{1/p}
+
(|a|+|b|)(C_\Gamma\delta)^{1/p}.
\]
\end{theorem}

\begin{proof}
The equal amplitude case follows from \cref{thm:scalar_region_formula} and the tube inclusion. For unequal amplitudes, split the error into its value mismatch on \(E\cap F\) and the two disagreement pieces. Apply the triangle inequality and the same tube estimate.
\end{proof}

\begin{example}[A perturbed disk]\label{ex:perturbed_disk}
Let \(E\) and \(F\) be concentric disks of radii \(r\) and \(r+\delta\) in \(\RR^2\), with \(\delta>0\). Then
\[
\mu(E\mathbin\triangle F)
=
\pi\bigl((r+\delta)^2-r^2\bigr)
=
2\pi r\delta+\pi\delta^2.
\]
For equal amplitude \(a\), the exact \(L^2\) packet error is
\[
|a|\sqrt{2\pi r\delta+\pi\delta^2}.
\]
This gives the expected square root dependence on boundary displacement.
\end{example}

\section{Stable recovery of separated jumps}\label{extension:section:09}\label{sec:jump_recovery}

\begin{keyidea}[title={The detector needs a visible interval of thresholds}]
A smooth cell and bounded sample noise give one upper contrast bound. A jump cell gives a lower contrast bound after the same uncertainty is subtracted. Recovery is possible when the lower jump bound exceeds the upper smooth bound. The gap is a condition on the observation model, rather than an automatic consequence of packet coordinates.
\end{keyidea}
\begin{figure}[htbp]\centering
\includegraphics[width=.93\linewidth]{monograph_figures/jump_contrast_gap.pdf}
\caption{A sufficient threshold interval from exact rational contrast bounds. With the displayed model parameters, smooth and noise contrasts are at most $1/10$ and jump contrasts are at least $9/10$. Every threshold strictly between those values separates the cells under the sampling assumptions stated below. The figure does not assert an optimal threshold for other observation models.}\label{monograph:fig:contrast}
\end{figure}


The previous sections describe the error after a packet has been estimated. We now give a direct recovery theorem for one dimensional separated jumps. The detector is intentionally simple because the aim is to expose the exact contrast gap.

Let
\[
t_k=kh,
\qquad
k=0,\ldots,n,
\qquad
h=\frac1n.
\]
Assume
\[
f(t)
=
s(t)+\sum_{q=1}^{Q}a_qH_{\tau_q}(t),
\qquad
H_\tau(t)=\one_{[\tau,1]}(t),
\]
where \(s\) is Lipschitz with constant \(M\), and \(Q\geq1\) is finite with nonzero jump amplitudes. The zero-jump case requires no minimum-amplitude condition. Suppose no jump lies at a sample point and no grid cell contains more than one jump. Observe
\[
y_k=f(t_k)+\eta_k.
\]
Define the adjacent contrast
\[
D_k=y_{k+1}-y_k.
\]

\begin{lemma}[Contrast gap]\label{lem:contrast_gap}
Assume \(|\eta_k|\leq\varepsilon\) and set
\[
b_h=Mh+2\varepsilon.
\]
If the cell \((t_k,t_{k+1})\) contains no jump, then
\[
|D_k|\leq b_h.
\]
If it contains one jump of amplitude \(a_q\), then
\[
|D_k-a_q|\leq b_h
\]
and therefore
\[
|D_k|\geq |a_q|-b_h.
\]
\end{lemma}

\begin{proof}
In a cell without a jump,
\[
D_k=s(t_{k+1})-s(t_k)+\eta_{k+1}-\eta_k.
\]
The Lipschitz and noise bounds give \(|D_k|\leq Mh+2\varepsilon\). In a cell with one jump, the same formula has the additional term \(a_q\), which proves the second estimate.
\end{proof}

\begin{theorem}[Deterministic jump cell recovery]\label{thm:deterministic_jump_recovery}
Let
\[
a_{\min}=\min_q|a_q|
\]
and assume
\[
a_{\min}>2b_h.
\]
Choose a threshold \(\theta\) satisfying
\[
b_h<\theta<a_{\min}-b_h.
\]
Then
\[
\widehat\Kcal
=
\{k:|D_k|>\theta\}
\]
is exactly the set of grid cells containing jumps. If \(\widehat\tau_q\) is the midpoint of the detected cell and \(\widehat a_q=D_k\), then after matching by order,
\[
|\widehat\tau_q-\tau_q|
\leq
\frac h2
\]
and
\[
|\widehat a_q-a_q|
\leq
b_h.
\]
\end{theorem}

\begin{proof}
By \cref{lem:contrast_gap}, every cell without a jump has contrast at most \(b_h<\theta\). Every cell with a jump has contrast at least \(a_{\min}-b_h>\theta\). Hence the detected set is exact. The midpoint of a cell is within \(h/2\) of every point in that cell. The amplitude estimate is the second inequality in \cref{lem:contrast_gap}.
\end{proof}

\begin{corollary}[Jump count and sign recovery]\label{cor:jump_count_sign}
Under the hypotheses of \cref{thm:deterministic_jump_recovery}, the number of detected cells equals \(Q\), and
\[
\sgn(\widehat a_q)=\sgn(a_q)
\]
for every real valued jump.
\end{corollary}

\begin{proof}
Exact cell recovery gives the jump count. Since
\[
|\widehat a_q-a_q|\leq b_h<|a_q|,
\]
the amplitude estimate cannot cross zero.
\end{proof}

\begin{corollary}[Packet reconstruction after detection]\label{cor:reconstruction_after_detection}
Let
\[
j=\sum_{q=1}^{Q}a_qH_{\tau_q},
\qquad
\widehat j=\sum_{q=1}^{Q}\widehat a_qH_{\widehat\tau_q}.
\]
Under the hypotheses of \cref{thm:deterministic_jump_recovery}, for \(1\leq p<\infty\),
\[
\|j-\widehat j\|_p
\leq
Qb_h
+
\left(
\frac h2\sum_{q=1}^{Q}|a_q|^p
\right)^{1/p}.
\]
If \(Q_m\) is bounded and linear and
\[
\widehat f_m=Q_m(f-\widehat j)+\widehat j,
\]
then
\[
\begin{aligned}
\|f-\widehat f_m\|_p
\leq{}&
\|s-Q_ms\|_p\\
&+(1+\|Q_m\|)
\left[
Qb_h+
\left(
\frac h2\sum_q|a_q|^p
\right)^{1/p}
\right].
\end{aligned}
\]
\end{corollary}

\begin{proof}
Write
\[
j-\widehat j
=
\sum_q(a_q-\widehat a_q)H_{\widehat\tau_q}
+
\sum_q a_q(H_{\tau_q}-H_{\widehat\tau_q}).
\]
The first sum has norm at most \(\sum_q|a_q-\widehat a_q|\leq Qb_h\). The location error interval for the \(q\)th jump is contained in its detected grid cell and has length at most \(h/2\). Distinct detected cells are disjoint, so the location errors have disjoint supports. Their \(p\)th powers therefore add exactly, which gives the displayed location term. The reconstruction estimate follows from \cref{thm:estimated_packet_reconstruction}.
\end{proof}

\begin{remark}[Why finite \(L^p\) matters]\label{rem:finite_p_recovery}
A shifted step does not converge in the supremum norm until the location is exact. The finite \(L^p\) error does converge because the disagreement is confined to an interval whose length tends to zero.
\end{remark}

\begin{corollary}[Deterministic consistency]\label{cor:deterministic_consistency}
Assume \(h\to0\) along grids satisfying the preceding sampling assumptions, \(\varepsilon_h\to0\), the number of jumps is fixed, and the nonzero jump amplitudes and positive minimum separation remain fixed. Then the detected packet converges to the true packet in every finite \(L^p\) norm. If \(Q_m\) has a vanishing smooth approximation error and uniformly bounded norm, the complete reconstruction also converges in \(L^p\).
\end{corollary}

\begin{proof}
The gap condition holds for sufficiently small \(h\). Then \(b_h\to0\), and the bound in \cref{cor:reconstruction_after_detection} tends to zero.
\end{proof}

\begin{figure}
\centering
\includegraphics[width=0.83\textwidth]{09_jump_detection.pdf}
\caption{Separated jumps create isolated adjacent sample contrasts above the smooth and noise background.}
\label{fig:jump_detection}
\end{figure}

\begin{example}[A numerical threshold certificate]\label{ex:jump_threshold_certificate}
Take \(h=1/256\), \(M=0.5\), and \(\varepsilon=0.002\). Then
\[
b_h
=\frac{0.5}{256}+0.004
<0.00596.
\]
If every jump has magnitude at least \(0.08\), then \(0.08>2b_h\). Any threshold between \(0.00596\) and \(0.07404\) recovers the exact jump cells. The location error is at most
\[
\frac1{512}
\]
and every amplitude error is less than \(0.00596\).
\end{example}

\subsection{Sub Gaussian noise}

Assume the noise variables are independent, centered, and \(\sigma\) sub Gaussian:
\[
\mathbb E\e^{\lambda\eta_k}
\leq
\e^{\sigma^2\lambda^2/2}.
\]

\begin{theorem}[Probability recovery bound]\label{thm:probability_recovery}
Fix \(0<\delta<1\) and define
\[
t_{n,\delta}
=
2\sigma
\sqrt{\log\left(\frac{2n}{\delta}\right)}
\]
and
\[
b_{h,\delta}=Mh+t_{n,\delta}.
\]
With probability at least \(1-\delta\), every contrast satisfies the conclusions of \cref{lem:contrast_gap} with \(b_h\) replaced by \(b_{h,\delta}\).

If
\[
a_{\min}>2b_{h,\delta},
\]
then thresholding between \(b_{h,\delta}\) and \(a_{\min}-b_{h,\delta}\) recovers every jump cell with probability at least \(1-\delta\). The same location and amplitude bounds hold with \(b_{h,\delta}\).
\end{theorem}

\begin{proof}
If \(\sigma=0\), every centered zero-sub-Gaussian noise variable vanishes almost surely, and the noise-free deterministic conclusion applies. For \(\sigma>0\), the difference \(\eta_{k+1}-\eta_k\) is sub Gaussian with variance proxy \(2\sigma^2\), so
\[
\mathbb P\bigl(|\eta_{k+1}-\eta_k|>t\bigr)
\leq
2\exp\left(-\frac{t^2}{4\sigma^2}\right).
\]
A union bound over the \(n\) cells gives
\[
\mathbb P\left(
\max_k|\eta_{k+1}-\eta_k|>t_{n,\delta}
\right)
\leq
\delta.
\]
On the complementary event, repeat the deterministic proof.
\end{proof}

\begin{corollary}[Repeated measurements]\label{cor:repeated_measurements}
Suppose \(R\) observations are averaged at every sample location, with all noise variables independent across sites and repetitions and each centered and \(\sigma\) sub Gaussian. Then \cref{thm:probability_recovery} holds with \(\sigma\) replaced by \(\sigma/\sqrt R\). In particular, a sufficient recovery condition is
\[
a_{\min}
>
2Mh+
\frac{4\sigma}{\sqrt R}
\sqrt{\log\left(\frac{2n}{\delta}\right)}.
\]
\end{corollary}

\begin{proof}
The average of \(R\) independent \(\sigma\) sub Gaussian variables is \(\sigma/\sqrt R\) sub Gaussian. Substitute this parameter in \cref{thm:probability_recovery}.
\end{proof}

\begin{remark}\label{rem:noise_scaling}
The displayed uniform noise certificate scales as \(\sigma\sqrt{\log n}\). At fixed positive \(\sigma\), it does not imply consistency of this detector. Decreasing the noise scale or averaging repeated measurements gives sufficient regimes; other averaging schemes require their own analysis.
\end{remark}

\begin{corollary}[Probability consistency]\label{cor:probability_consistency}
Let \(\sigma_n\sqrt{\log n}\to0\) and \(h=1/n\). For every fixed \(0<\delta<1\), along grids satisfying the preceding sampling assumptions, all jump cells are recovered with probability at least \(1-\delta\) for sufficiently large \(n\). For the fixed finite jump model above, along grids satisfying the preceding sampling assumptions, the packet error converges to zero in probability in every finite \(L^p\) norm.
\end{corollary}

\begin{proof}
For each fixed \(\delta>0\), \(b_{h,\delta}\to0\). Apply \cref{thm:probability_recovery,cor:reconstruction_after_detection}. For any error tolerance \(\epsilon>0\), the high-probability bound is eventually smaller than \(\epsilon\), so the limit superior of the probability that the packet error exceeds \(\epsilon\) is at most \(\delta\). Since \(\delta\) is arbitrary, convergence is in probability.
\end{proof}

\subsection{Signals with values in a normed space}

The recovery proof does not depend on an ordering of the signal values.

\begin{proposition}[Vector jump recovery]\label{prop:vector_jump_recovery}
Let the values lie in a normed space \(V\). Assume
\[
\|s(t)-s(u)\|_V\leq M|t-u|,
\qquad
\|\eta_k\|_V\leq\varepsilon,
\]
and let the jump vectors satisfy \(\|a_q\|_V\geq a_{\min}\). Thresholding \(\|D_k\|_V\) gives the same recovery theorem under the condition
\[
a_{\min}>2(Mh+2\varepsilon).
\]
\end{proposition}

\begin{proof}
Only the triangle inequality was used in \cref{lem:contrast_gap,thm:deterministic_jump_recovery}.
\end{proof}

\section{Haar localization and annihilating differences}\label{extension:section:10}\label{sec:localization}

The adjacent difference detector is local in the sample domain. Wavelet coefficients and higher differences give related multiscale packets. In the Haar subsections the signal is defined on \(\RR\). Here the step convention is \(H_\tau(t)=\one_{[\tau,\infty)}(t)\), whose restriction to \((0,1)\) agrees with the preceding interval step. A chosen extension beyond the observation interval must use this half-line step convention, and the smooth extension must satisfy the stated Lipschitz hypothesis on every displayed wavelet support.

\subsection{Haar coefficient separation}

Let
\[
I_{j,k}
=
[k2^{-j},(k+1)2^{-j})
\]
and let \(\psi_{j,k}\) be the normalized Haar wavelet supported on \(I_{j,k}\), positive on the left half and negative on the right half.

\begin{theorem}[Smooth Haar coefficient bound]\label{thm:smooth_haar_bound}
If \(s\) is Lipschitz with constant \(M\), then
\[
|\langle s,\psi_{j,k}\rangle|
\leq
\frac M4\,2^{-3j/2}.
\]
\end{theorem}

\begin{proof}
Let \(L=2^{-j}\) and let \(a\) be the left endpoint of the support. Pair every point \(t\in[a,a+L/2]\) with \(t+L/2\). Then
\[
\langle s,\psi_{j,k}\rangle
=
L^{-1/2}
\int_a^{a+L/2}
\overline{s(t)-s(t+L/2)}\,\dd t.
\]
The integrand has magnitude at most \(ML/2\), and the integration interval has length \(L/2\). The result is \(ML^{3/2}/4\).
\end{proof}

\begin{theorem}[Exact Haar coefficient of a step]\label{thm:haar_step_coefficient}
Let \(\tau\in I_{j,k}\), let \(L=2^{-j}\), and write
\[
u=\frac{\tau-kL}{L}\in[0,1].
\]
Then
\[
\left|
\left\langle
 aH_\tau,\psi_{j,k}
\right\rangle
\right|
=
|a|L^{1/2}\min\{u,1-u\}.
\]
Every Haar coefficient whose support does not cross \(\tau\) is zero.
\end{theorem}

\begin{proof}
Away from the jump, the step is constant on the support and the Haar wavelet has zero mean. When the jump lies in the support, integrate the positive and negative halves directly. If \(u\leq1/2\), the magnitude is \(|a|uL^{1/2}\). If \(u\geq1/2\), it is \(|a|(1-u)L^{1/2}\).
\end{proof}

\begin{corollary}[Scale separation]\label{cor:haar_scale_separation}
Suppose \(f=s+aH_\tau+\eta\), where \(s\) is \(M\) Lipschitz and \(\|\eta\|_\infty\leq\varepsilon\). If the jump lies at relative distance at least \(\rho\) from both ends of \(I_{j,k}\), then
\[
|\langle f,\psi_{j,k}\rangle|
\geq
\left(
|a|\rho
-\varepsilon
-\frac M4L
\right)L^{1/2}.
\]
A wavelet support without a jump satisfies
\[
|\langle f,\psi_{j,k}\rangle|
\leq
\left(
\varepsilon+\frac M4L
\right)L^{1/2}.
\]
Thus a threshold gap exists whenever
\[
|a|\rho
>
2\varepsilon+\frac M2L.
\]
\end{corollary}

\begin{proof}
Combine \cref{thm:smooth_haar_bound,thm:haar_step_coefficient}. The noise coefficient is at most
\[
\|\eta\|_\infty\|\psi_{j,k}\|_1
=
\varepsilon L^{1/2}.
\]
The final inequality makes the lower bound on a jump support larger than the upper bound on a smooth support.
\end{proof}

\begin{figure}
\centering
\includegraphics[width=0.68\textwidth]{10_haar_scale_separation.pdf}
\caption{The smooth contribution decays like \(L^{3/2}\), while a separated jump contribution decays like \(L^{1/2}\).}
\label{fig:haar_scale_separation}
\end{figure}

\subsection{Two shifted grids}

A jump on a dyadic boundary can be invisible to a wavelet at that scale. One additional grid removes this exact alignment.

\begin{theorem}[Half cell shift guarantee]\label{thm:half_shift_guarantee}
Fix a cell length \(L>0\). Consider the grid with boundaries \(kL\) and the grid with boundaries \((k+1/2)L\). Every point \(\tau\in\RR\) lies in a cell of at least one grid at distance at least \(L/4\) from both cell boundaries.
\end{theorem}

\begin{proof}
Reduce \(\tau\) modulo \(L\) and write its representative as \(x\in[0,L)\). If
\[
L/4\leq x\leq 3L/4,
\]
then the distance from \(x\) to both boundaries of its cell in the original grid is at least \(L/4\). If \(x<L/4\) or \(x>3L/4\), then \(x\) lies in the cell of the shifted grid bounded by \(-L/2\) and \(L/2\), or in its translate, at distance at least \(L/4\) from both shifted boundaries. One of the two grids therefore has the required property.
\end{proof}

\begin{corollary}[Alignment free two grid Haar detection]\label{cor:two_grid_haar_detection}
Assume the support length \(L\) is smaller than the minimum jump separation. Analyze the signal with Haar wavelets on the original grid and on the half shifted grid. Every jump appears in a support with \(\rho\geq1/4\) in one of the two systems. Therefore a threshold gap exists independently of alignment whenever
\[
\frac{|a|}{4}
>
2\varepsilon+\frac M2L.
\]
\end{corollary}

\begin{proof}
Apply \cref{thm:half_shift_guarantee} and then \cref{cor:haar_scale_separation} with \(\rho=1/4\). The separation assumption prevents a selected support from containing two jumps.
\end{proof}

\begin{figure}
\centering
\includegraphics[width=0.78\textwidth]{11_shifted_grids.pdf}
\caption{The boundaries of the two grids interlace. Every jump is well inside a cell in at least one grid.}
\label{fig:shifted_grids}
\end{figure}

\subsection{Higher differences}

Let
\[
(\Delta x)(n)=x(n+1)-x(n).
\]
Then
\[
\Delta^m x(n)
=
\sum_{r=0}^{m}
(-1)^{m-r}\binom mr x(n+r).
\]

\begin{theorem}[Polynomial annihilation and boundary support]\label{thm:polynomial_annihilation}
Let \(x\) be a sequence that agrees with a polynomial of degree at most \(d\) on each interval of a finite partition of \(\ZZ\). Then
\[
\Delta^{d+1}x(n)=0
\]
whenever the complete stencil
\[
\{n,n+1,\ldots,n+d+1\}
\]
lies inside one polynomial piece. Hence \(\Delta^{d+1}x\) is supported in a band of at most \(d+1\) addresses before each breakpoint.
\end{theorem}

\begin{proof}
The \((d+1)\)st forward difference of a polynomial of degree at most \(d\) is zero. The only stencils for which this argument does not apply are those crossing a breakpoint.
\end{proof}

\begin{proposition}[Noise amplification by differences]\label{prop:difference_noise}
If \(\|\eta\|_\infty\leq\varepsilon\), then
\[
\|\Delta^m\eta\|_\infty
\leq
2^m\varepsilon.
\]
Therefore no false response above \(2^{d+1}\varepsilon\) can occur on a stencil lying inside one polynomial piece.
\end{proposition}

\begin{proof}
The absolute values of the binomial coefficients in \(\Delta^m\) sum to \(2^m\).
\end{proof}

\begin{corollary}[Robust breakpoint bands]\label{cor:robust_breakpoint_bands}
Assume every true breakpoint produces at least one response of magnitude \(\rho\) in its \((d+1)\)st difference band. If
\[
\rho>2^{d+2}\varepsilon,
\]
then thresholding at any level between \(2^{d+1}\varepsilon\) and \(\rho-2^{d+1}\varepsilon\) detects a nonempty cluster near every breakpoint and creates no cluster away from the breakpoint bands.
\end{corollary}

\begin{proof}
Use \cref{thm:polynomial_annihilation,prop:difference_noise}. Noise cannot cross the threshold away from a breakpoint. At a response of magnitude \(\rho\), the perturbed magnitude remains at least \(\rho-2^{d+1}\varepsilon\), which is above the threshold.
\end{proof}

\section{Packet separation and the Gibbs phenomenon}\label{extension:section:11}\label{sec:gibbs}

A Fourier partial sum is global. When the target has a jump, the partial sum must approximate that jump using continuous trigonometric polynomials. This creates the classical Gibbs oscillation. Standard accounts and recovery methods appear in \cite{Zygmund2002,GottliebShu1997}.

\begin{standardfact}[Classical Gibbs behavior]\label{sf:gibbs_behavior}
For a periodic piecewise \(C^1\) function with an isolated jump of size \(a\), ordinary Fourier partial sums have an overshoot that approaches a fixed nonzero multiple of \(|a|\). The width of the oscillatory region shrinks, but the maximum overshoot does not tend to zero.
\end{standardfact}

The retained packet identity changes the approximation problem because the discontinuity is no longer represented by the Fourier partial sum.

\begin{theorem}[Uniform convergence with exact jump packets]\label{thm:uniform_packet_convergence}
Let \(f=s+j\) on a compact domain, where \(j\) is any bounded singular packet. Suppose \(Q_ms\to s\) uniformly. Then
\[
\widetilde f_m=Q_ms+j
\]
converges uniformly to \(f\), even when \(f\) is discontinuous:
\[
\|f-\widetilde f_m\|_\infty
=
\|s-Q_ms\|_\infty
\longrightarrow0.
\]
There is no error proportional to the jump size because the jump is shared exactly by \(f\) and \(\widetilde f_m\).
\end{theorem}

\begin{proof}
Apply \cref{thm:retained_packet_identity} in the supremum norm.
\end{proof}

\begin{theorem}[Estimated packet Gibbs bound]\label{thm:estimated_packet_gibbs}
Let \(Q_m\) be linear and uniformly bounded on \(L^p\), with
\[
\sup_m\|Q_m\|\leq C_Q.
\]
For \(f=s+j\) and an estimated packet \(\widehat j_m\), define
\[
\widehat f_m
=Q_m(f-\widehat j_m)+\widehat j_m.
\]
Then
\[
\|f-\widehat f_m\|_p
\leq
\|s-Q_ms\|_p
+(1+C_Q)\|j-\widehat j_m\|_p.
\]
Consequently, smooth approximation and packet recovery may be improved independently. No claim of Gibbs removal is justified unless both terms are controlled in the norm under consideration.
\end{theorem}

\begin{proof}
This is \cref{thm:estimated_packet_reconstruction} with the uniform operator bound.
\end{proof}

\begin{corollary}[Boundary band error]\label{cor:boundary_band_error}
Let \(\Gamma\) be the singular boundary and assume
\[
\mu(N_\delta(\Gamma))\leq C_\Gamma\delta
\]
for small \(\delta\). For finite \(p\),
\[
\|f-\widetilde f_m\|_{L^p(N_\delta(\Gamma))}
\leq
(C_\Gamma\delta)^{1/p}
\|s-Q_ms\|_\infty.
\]
\end{corollary}

\begin{proof}
The error equals \(s-Q_ms\). Bound it by its supremum and multiply by the measure of the boundary tube.
\end{proof}

\begin{example}[Square wave Fourier tail]\label{ex:square_wave_tail}
Let
\[
j(t)=\one_{[-a,a]}(t)
\]
on the circle and let \(s\) be smooth. If \(Q_m\) is a Fourier partial sum, then \(Q_m(s+j)\) has the ordinary oscillatory boundary layer. In contrast,
\[
Q_ms+j
\]
has error exactly \(s-Q_ms\). The visual comparison in \cref{fig:gibbs_comparison} illustrates this change in approximation target. It is not a new Fourier summability theorem.
\end{example}

\begin{figure}
\centering
\includegraphics[width=0.82\textwidth]{12_gibbs_comparison.pdf}
\caption{Global Fourier truncation represents the jump and exhibits an oscillatory boundary layer. Packet separation retains the jump and approximates only the smooth component.}
\label{fig:gibbs_comparison}
\end{figure}

\section{Nonlinear interactions in residue coordinates}\label{extension:section:12}\label{sec:nonlinear}

Linear transfer matrices do not directly describe nonlinear products. Residue coordinates still obey exact conservation laws because integer addresses add modulo the packet modulus.

\begin{theorem}[Residue law for Cauchy products]\label{thm:cauchy_residue_law}
Let
\[
F(z)=\sum_{n\geq0}f_nz^n,
\qquad
G(z)=\sum_{n\geq0}g_nz^n
\]
be formal power series. The same statement holds for power series that converge on a common disk. Decompose both series into residue packets modulo \(N\):
\[
F=\sum_{r=0}^{N-1}F_r,
\qquad
G=\sum_{s=0}^{N-1}G_s.
\]
If \(H=FG\), then its residue packet \(H_t\) is
\[
\boxed{
H_t
=
\sum_{r+s\equiv t\!\!\pmod N}F_rG_s.
}
\]
\end{theorem}

\begin{proof}
Every monomial in \(F_rG_s\) has degree congruent to \(r+s\) modulo \(N\). Conversely, every coefficient of degree congruent to \(t\) is formed from pairs whose residues sum to \(t\). Grouping the Cauchy product by this residue gives the formula.
\end{proof}

\begin{theorem}[Shifted pointwise monomial law]\label{thm:shifted_monomial_law}
Let
\[
y(n)=\prod_{j=1}^{m}x(n-a_j)
\]
and write
\[
a_j=Nc_j+b_j,
\qquad
0\leq b_j<N.
\]
For output residue \(r\), define
\[
r_j\equiv r-b_j\pmod N,
\qquad
0\leq r_j<N,
\]
and let
\[
\kappa_{r,b_j}
=
\begin{cases}
1,&r<b_j,\\
0,&r\geq b_j.
\end{cases}
\]
Then the output channel is
\[
\boxed{
y_r(q)
=
\prod_{j=1}^{m}
x_{r_j}
\bigl(q-c_j-\kappa_{r,b_j}\bigr).
}
\]
\end{theorem}

\begin{proof}
For \(n=Nq+r\),
\[
n-a_j
=N(q-c_j)+r-b_j.
\]
If \(r\geq b_j\), this address has residue \(r-b_j\) and coarse index \(q-c_j\). If \(r<b_j\), add \(N\) to the residue and subtract one from the coarse index. Multiplying the resulting channel values gives the formula.
\end{proof}

\begin{corollary}[Polynomial shift filters]\label{cor:polynomial_shift_filters}
Every finite polynomial in shifted samples has an exact finite expression in residue coordinates. Each monomial selects a deterministic list of input residues and coarse shifts, and the complete output is their finite sum.
\end{corollary}

\begin{proof}
Apply \cref{thm:shifted_monomial_law} to each monomial and add the results.
\end{proof}

\begin{example}[Quadratic nearest neighbor interaction]\label{ex:quadratic_neighbor}
For
\[
y(n)=x(n)x(n-1)
\]
and modulus \(N\),
\[
y_r(q)
=
\begin{cases}
x_r(q)x_{r-1}(q),&r\geq1,\\
x_0(q)x_{N-1}(q-1),&r=0.
\end{cases}
\]
The carry at residue zero is the nonlinear analogue of a transfer crossing a coarse cell boundary.
\end{example}

\section{Finite numerical illustrations}\label{extension:section:13}\label{sec:numerical}

The figures and tables in this section are finite checks and examples. They are not substitutes for the preceding proofs. Their purpose is to show the scale and shape of the terms appearing in the exact estimates.

A fixed depth Haar packet transform and a Fourier basis were applied to the same piecewise smooth test signal. The largest coefficients were retained in each representation.

\begin{figure}
\centering
\includegraphics[width=0.70\textwidth]{13_sparse_comparison.pdf}
\caption{One finite sparse approximation comparison. Fourier coordinates are better at the smallest budget, while the fixed Haar packet basis becomes better after localized edge coordinates are resolved.}
\label{fig:sparse_comparison}
\end{figure}

\begin{table}
\centering
\caption{Finite sparse approximation record for one test signal.}
\label{tab:sparse_record}
\begin{tabular}{rrr}
\toprule
retained & Haar packet RMSE & Fourier RMSE\\
\midrule
16 & \(1.7999\times10^{-1}\) & \(6.6522\times10^{-2}\)\\
32 & \(3.8414\times10^{-2}\) & \(4.9153\times10^{-2}\)\\
64 & \(1.6267\times10^{-2}\) & \(3.8056\times10^{-2}\)\\
128 & \(1.7387\times10^{-3}\) & \(2.5631\times10^{-2}\)\\
256 & \(1.5870\times10^{-4}\) & \(1.7424\times10^{-2}\)\\
\bottomrule
\end{tabular}
\end{table}

The table does not prove that Haar packets dominate Fourier coordinates for every signal. Smooth periodic oscillations may favor Fourier coordinates. Curved edges in two dimensions may favor directional systems. The geometric approximation results in \cite{Donoho1999,CandesDonoho2004} give much deeper statements for specific image classes.

\begin{table}
\centering
\caption{Finite edge reconstruction record from the original computational study.}
\label{tab:edge_record}
\small
\begin{tabular}{lccc}
\toprule
case & global & exact packet & estimated packet\\
\midrule
one dimensional step & \(5.4023\times10^{-2}\) & \(1.3692\times10^{-4}\) & \(2.4029\times10^{-2}\)\\
two dimensional disk & \(3.7643\times10^{-2}\) & \(6.4752\times10^{-5}\) & \(3.4043\times10^{-3}\)\\
multiple regions & \(3.3740\times10^{-2}\) & \(5.2118\times10^{-7}\) & not evaluated\\
\bottomrule
\end{tabular}
\end{table}

The estimated packet rows used a narrow positive contrast synthetic model. They are not evidence for a general segmentation theorem. The exact packet rows verify \cref{thm:retained_packet_identity} numerically after the packet is supplied.

\section{Cyclic Cayley coordinates}\label{extension:section:14}\label{sec:cayley}

\begin{interpretation}[title={Keep the rotation and its chart separate}]
For the phase factorization below, $M$ and $N$ are positive phase orders. The reflection coordinate rotates in the unit disk. The Cayley map expresses that same motion as a fractional linear impedance action. A pole of the chart does not mean that the projective action fails. This realization is independent of the separated jump detector.
\end{interpretation}

This section records a secondary realization from the original signal manuscript. It is logically independent of the recovery theorems. The classical Smith chart uses the Cayley transform between normalized impedance and a reflection coordinate \cite{Smith1969}.

Let
\[
\mathbb H
=
\{z\in\CC:\Re z>0\}
\]
and define
\[
C(z)=\frac{z-1}{z+1}.
\]
It maps \(\mathbb H\) conformally onto the unit disk \(\mathbb D\), with inverse
\[
C^{-1}(w)=\frac{1+w}{1-w}.
\]
For \(\theta\in\RR\), let
\[
R_\theta(w)=\e^{\ii\theta}w.
\]

\begin{theorem}[Cyclic Mobius action]\label{thm:cyclic_mobius}
The pullback
\[
\Phi_\theta
=C^{-1}R_\theta C
\]
is the fractional linear map
\[
\boxed{
\Phi_\theta(z)
=
\frac{(1+\e^{\ii\theta})z+(1-\e^{\ii\theta})}
{(1-\e^{\ii\theta})z+(1+\e^{\ii\theta})}.
}
\]
It preserves \(\mathbb H\) and satisfies
\[
\Phi_\theta\Phi_\varphi
=
\Phi_{\theta+\varphi}.
\]
For an integer \(N\geq1\) and \(\theta=2\pi/N\), the map has order \(N\).
\end{theorem}

\begin{proof}
Substitute \(\e^{\ii\theta}(z-1)/(z+1)\) into \(C^{-1}\) and simplify. Preservation of \(\mathbb H\) follows because \(C\) maps it to \(\mathbb D\), rotation preserves \(\mathbb D\), and \(C^{-1}\) maps back. The group law follows by conjugating the rotation group law. When \(\theta=2\pi/N\), the \(N\)th power is the identity. Conversely, if its \(k\)th power is the identity, conjugation by \(C\) and evaluation at \(w=1/2\) give \(\e^{2\pi\ii k/N}=1\), so \(N\) divides \(k\). Hence the order is exactly \(N\).
\end{proof}

\begin{corollary}[Impedance and admittance]\label{cor:impedance_admittance}
For \(N=2\), \(\theta=\pi\) and
\[
\Phi_\pi(z)=\frac1z.
\]
Thus the two phase action is normalized impedance and admittance inversion.
\end{corollary}

\begin{proof}
Insert \(\e^{\ii\pi}=-1\) in the formula of \cref{thm:cyclic_mobius}.
\end{proof}

\begin{proposition}[Radial and phase coordinates]\label{prop:radial_phase}
Write
\[
C(z)=\rho\e^{\ii\varphi},
\qquad
0\leq\rho<1.
\]
Then
\[
C(\Phi_\theta(z))
=
\rho\e^{\ii(\varphi+\theta)}.
\]
The radial mismatch \(\rho\) is invariant and the phase coordinate advances by \(\theta\).
\end{proposition}

\begin{proof}
This is the defining conjugacy \(C\Phi_\theta=R_\theta C\).
\end{proof}

\begin{proposition}[Coprime phase factorization]\label{prop:phase_factorization}
Let \(M\) and \(N\) be coprime. The cyclic phase labels modulo \(MN\) are equivalent to pairs of labels modulo \(M\) and modulo \(N\):
\[
\ZZ/(MN)\ZZ
\cong
\ZZ/M\ZZ\times\ZZ/N\ZZ.
\]
Under this identification, the rotation by \(2\pi/(MN)\) is determined by its two projected cyclic actions.
\end{proposition}

\begin{proof}
This is the Chinese remainder theorem applied to the additive phase label. A label modulo \(MN\) is uniquely determined by its two residues, and addition of labels is preserved.
\end{proof}

\begin{figure}
\centering
\includegraphics[width=0.61\textwidth]{14_smith_coordinates.pdf}
\caption{The cyclic action is an ordinary phase rotation in the reflection disk and a fractional linear action in the impedance half plane.}
\label{fig:smith_coordinates}
\end{figure}

\section{Limits and further directions}\label{extension:section:15}\label{sec:further_questions}

The one dimensional recovery theorem identifies an explicit admissible threshold interval for an adjacent contrast detector. More general recovery requires a broader observation model. A multiscale average can reduce fixed sample noise, but its bias depends on the smoothness class and the separation of nearby jumps. A useful next theorem should select the averaging scale from the data while retaining a finite probability certificate.

Higher dimensional recovery requires geometry. The overlap and tube formulas in \cref{sec:packet_error} become quantitative once a boundary estimate is known. What remains is a stable method that recovers topology, orientation, curvature, and region values under a declared noise model. Wedgelets and curvelets show why directional geometry is essential for broad image classes \cite{Donoho1999,CandesDonoho2004}.

The best basis recursion is exact for additive objectives. Cross packet leakage, shared computation, and global hardware constraints introduce couplings between subtrees. Those problems require a larger compiler or an approximation theorem. The cost stability results in \cref{thm:best_basis_cost_stability,cor:best_basis_selection_stability} identify one condition under which the additive decision survives perturbation.

The higher order singular packet theorem removes finitely many trace jumps. A more complete approximation theory should connect packet norms to geometric measure quantities in \(\BV\), piecewise Sobolev spaces, and free discontinuity models. It should also determine when a redundant packet dictionary can identify the derivative order of a singularity rather than merely its location.

Finally, the two shifted Haar grids remove one exact alignment failure, but they are redundant. A complete redundant theory should provide stable synthesis, noise aware coefficient aggregation, and a best basis rule that does not count the same singularity several times.

\section{Conclusion}\label{extension:section:16}\label{sec:conclusion}

Finite index sublattices give exact coordinates for lattice signals. In those coordinates, convolution is a matrix of transfer kernels. Perfect reconstruction is a matrix inverse identity. Approximate reconstruction is controlled by a multiplier norm. Frame bounds are pointwise singular value bounds. Lifting steps are exact triangular coordinate changes. Every admissible packet tree gives a complete coordinate system, and every additive tree objective has an exact dynamic recursion.

The singular packet theory separates representation from recovery. A finite family of truncated powers records jumps in a function and its derivatives, leaving a globally Sobolev remainder under explicit trace assumptions. If the packet is known and retained, the error is exactly the smooth approximation error. If the packet is estimated, exact region overlap and location formulas quantify the remaining defect, while the reconstruction operator contributes a transparent stability factor.

For separated one dimensional jumps, a deterministic contrast gap gives exact cell recovery, midpoint location accuracy, and amplitude control. A sub Gaussian version follows from one union bound, and repeated measurements reduce the effective noise scale. Haar coefficients provide a second separation law because smooth coefficients decay one full power of scale faster than jump coefficients. Two shifted grids remove boundary alignment at a fixed scale. Higher differences annihilate piecewise polynomial interiors and confine their response to breakpoint bands.

The resulting chain is
\[
\begin{gathered}
\text{lattice coordinates}
\longrightarrow
\text{transfer blocks}
\longrightarrow
\text{stable filter banks}
\longrightarrow
\text{packet trees}
\\
\longrightarrow
\text{singular packet model}
\longrightarrow
\text{local recovery}
\longrightarrow
\text{smooth approximation rate}.
\end{gathered}
\]
Every arrow is exact or quantitatively controlled under the hypotheses stated in the corresponding result.


\clearpage
\part*{Part III: Forty worked visual plates}\addcontentsline{toc}{part}{Part III: Forty worked visual plates}
\begin{readerroadmap}[title={A second route through the mathematics}]
These forty plates develop the finite models used in the separate visual guide. The identifiers V1--V40 stay fixed across both documents. They lead from exact signal coordinates through filter transfers, adaptive decisions, packet errors, stable banks, and separated jump recovery. Every page gives the relevant manuscript locator. Finite arithmetic and diagrams explain the stated hypotheses and do not extend the formal contract.
\end{readerroadmap}
\clearpage
\begingroup
\setlength{\parskip}{3pt}
\setlength{\abovedisplayskip}{6pt}\setlength{\belowdisplayskip}{6pt}
\setstretch{1.01}
\fontsize{10}{12}\selectfont
\input{atlas-pages.tex}
\endgroup
\clearpage

\bibliographystyle{alpha}
\bibliography{references}


\end{document}
